\frfilename{mt254.tex}
\versiondate{23.2.16}
\copyrightdate{2002}

\def\chaptername{Ukuran produk}
\def\sectionname{Produk tak hingga}

\newsection{254}

Sekarang saya sampai pada gagasan dasar kedua dari
bab ini: uraian tentang ukuran produk pada produk suatu
keluarga ruang probabilitas (yang mungkin besar). Bagian ini dimulai dengan
suatu konstruksi yang sejalan dengan konstruksi dalam \S251 (254A-254F) dan
sifat penentunya dalam hal fungsi-fungsi \imp\ (254G). Saya membahas
ukuran biasa pada $\{0,1\}^I$
(254J-254K), ukuran subruang (254L), dan berbagai sifat
subproduk (254M-254T), termasuk kajian tentang operator ekspektasi
bersyarat yang terkait (254R-254T).

\leader{254A}{Definisi (a)} Misalkan $\familyiI{(X_i,\Sigma_i,\mu_i)}$ suatu
keluarga ruang
probabilitas. Tetapkan $X=\prod_{i\in I}X_i$, yaitu keluarga fungsi $x$ dengan
domain $I$ sedemikian sehingga $x(i)\in X_i$ untuk setiap $i\in I$. Dalam
konteks ini, saya akan mengatakan bahwa suatu {\bf silinder terukur} adalah subhimpunan $X$
yang dapat dinyatakan dalam bentuk

\Centerline{$C=\prod_{i\in I}C_i,$}

\noindent dengan $C_i\in\Sigma_i$ untuk setiap $i\in I$ dan
$\{i:C_i\ne X_i\}$ berhingga. Perhatikan bahwa untuk $C\subseteq X$ yang
tidak kosong, representasi ini unik. \prooflet{\Prf\ Andaikan
$C=\prod_{i\in I}C_i=\prod_{i\in I}C'_i$. Untuk setiap $i\in I$, tetapkan

\Centerline{$D_i=\{x(i):x\in C\}.$}

\noindent Tentu saja $D_i\subseteq C_i$. Karena $C\ne\emptyset$, kita
dapat memilih suatu $z\in C$. Jika $i\in I$ dan $t\in C_i$,
tinjau $x\in X$ yang didefinisikan dengan menetapkan

\Centerline{$x(i)=t$, \quad$x(j)=z(j)\text{ untuk }j\ne i;$}

\noindent maka $x\in C$, sehingga $t=x(i)\in D_i$. Jadi $D_i=C_i$ untuk
$i\in I$. Dengan cara yang sama, $D_i=C'_i$.\ \Qed}

\header{254Ab}{\bf (b)} Karena itu kita dapat mendefinisikan suatu fungsional
$\theta_0:\Cal C\to[0,1]$, dengan $\Cal C$ himpunan silinder
terukur, dengan menetapkan

\Centerline{$\theta_0C=\prod_{i\in I}\mu_iC_i$}

\noindent kapan pun $C_i\in\Sigma_i$ untuk setiap $i\in I$ dan
$\{i:C_i\ne X_i\}$ berhingga\cmmnt{,
dengan memperhatikan bahwa hanya berhingga banyak suku dalam produk yang berbeda dari $1$,
sehingga produk itu dapat diperlakukan dengan aman sebagai produk berhingga. Jika
$C=\emptyset$, salah satu $C_i$ harus kosong, sehingga $\theta_0C$ pastilah
$0$, meskipun representasi $C$ sebagai $\prod_{i\in I}C_i$
tidak lagi unik}.

\header{254Ac}{\bf (c)} Sekarang definisikan $\theta:\Cal PX\to[0,1]$ dengan

\Centerline{$\theta A
=\inf\{\sum_{n=0}^{\infty}\theta_0C_n:C_n\in\Cal C
  \text{ untuk setiap }n\in\Bbb N$, $A\subseteq\bigcup_{n\in\Bbb N}C_n\}$.}

\leader{254B}{Lema} Fungsional $\theta$ yang didefinisikan dalam 254Ac selalu
merupakan ukuran luar pada $X$.

\proof{ Gunakan tepat argumen yang sama dengan argumen dalam 251B di atas.}

\leader{254C}{Definisi} Misalkan
$\langle(X_i,\Sigma_i,\mu_i)\rangle_{i\in I}$ sebarang keluarga terindeks
ruang probabilitas, dan $X$ produk Kartesius
$\prod_{i\in I}X_i$. {\bf Ukuran produk} pada $X$ adalah ukuran yang
didefinisikan melalui metode \Caratheodory\cmmnt{ (113C)} dari
ukuran luar $\theta$ yang didefinisikan dalam 254A.

\cmmnt{
\leader{254D}{Catatan (a)} Dalam 254Ab, saya menegaskan bahwa jika $C\in\Cal C$ dan
tidak ada $C_i$ yang kosong, maka $C=\prod_{i\in I}C_i$ juga tidak
kosong. Inilah ‘Aksioma Pilihan’: produk dari
sebarang keluarga $\langle C_i\rangle_{i\in I}$ yang terdiri atas
himpunan-himpunan tidak kosong adalah tidak kosong, yaitu, terdapat suatu ‘fungsi pilihan’
$x$ dengan domain $I$ yang memilih satu anggota istimewa $x(i)$ dari setiap
$C_i$. Dalam jilid ini saya tidak berusaha untuk selalu teliti dalam
menunjukkan penggunaan aksioma pilihan. Bahkan, penggunaan di sini bukan
penggunaan yang mutlak penting; maksud saya, teori produk tak hingga, bahkan
produk tak terhitung, dari ruang probabilitas tidak berubah sifat
sepenuhnya tanpa aksioma pilihan penuh (asalkan, tentu saja,
kita mengizinkan diri menggunakan aksioma pilihan terhitung). Intinya
ialah bahwa yang sungguh-sungguh kita perlukan, dalam konteks sekarang, adalah bahwa
$X=\prod_{i\in I}X_i$ tidak kosong; dan dalam banyak konteks kita dapat
membuktikan hal ini, untuk kasus tertentu yang diperhatikan, tanpa menggunakan
aksioma pilihan, dengan benar-benar menunjukkan seorang anggota $X$. Kasus
paling sederhana yang menyulitkan adalah ketika $X_i$ merupakan subhimpunan Borel tak terkendali
dari $[0,1]$; dan bahkan ketika itu, jika disajikan dengan uraian
yang koheren, kita mungkin, dengan kerja yang memadai, dapat membangun seorang
anggota $X$. Namun, jelas bahwa proses semacam itu dapat memperlambat kita
secara berarti, dan untuk sementara saya kira tidak banyak manfaatnya menempuh
kesulitan sebesar itu.

\header{254Db}{\bf (b)}
Saya memberi bagian ini judul ‘produk tak hingga’, tetapi
gagasan-gagasannya berguna pula untuk $I$ berhingga; khususnya, saya perlu menyebutkan
kasus $\#(I)\le 2$.

\medskip

\quad {\bf (i)} Jika $I=\emptyset$, $X$ terdiri atas satu-satunya fungsi
dengan domain $I$, yaitu fungsi kosong. Jika kita mengidentifikasi fungsi dengan
grafnya, maka $X$ sesungguhnya adalah $\{\emptyset\}$; bagaimanapun, $X$
merupakan himpunan singleton, dengan $\lambda X=1$.

\medskip

\quad{\bf (ii)} Jika $I$ merupakan singleton $\{i\}$, maka kita dapat mengidentifikasi $X$
dengan $X_i$; $\Cal C$ menjadi teridentifikasi dengan $\Sigma_i$ dan $\theta_0$
dengan $\mu_i$, sehingga $\theta$ dapat diidentifikasi dengan $\mu_i^*$
dan ‘ukuran produk’ menjadi ukuran pada $X_i$
yang didefinisikan dari $\mu_i^*$, yaitu pelengkapan $\mu_i$\cmmnt{
(lihat 213Xa(iv))}.

\medskip

\quad{\bf (iii)} Jika $I$ merupakan himpunan beranggota dua $\{i,j\}$, maka kita dapat mengidentifikasi
$X$ dengan $X_i\times X_j$; dalam kasus ini definisi 254A dan 254C
persis sama dengan definisi
251A dan 251C, sehingga $\lambda$ di sini dapat diidentifikasi dengan
ukuran produk primitif sebagaimana didefinisikan dalam 251C. Karena $\mu_i$ dan
$\mu_j$ keduanya berhingga total, ukuran ini sama dengan ukuran produk c.l.d.\ yang
tercantum dalam 251F.
}%end of comment


\leader{254E}{Definisi} Misalkan $\langle X_i\rangle_{i\in I}$ sebarang
keluarga himpunan, dan $X=\prod_{i\in I}X_i$. Jika $\Sigma_i$ merupakan
subaljabar-$\sigma$ dari himpunan bagian $X_i$ untuk setiap $i\in I$, saya menuliskan
$\Tensorhat_{i\in I}\Sigma_i$ untuk aljabar-$\sigma$ himpunan bagian $X$
yang dibangkitkan oleh

\Centerline{$\{\{x:x\in X$, $x(i)\in E\}: i\in I$, $E\in\Sigma_i\}$.}

\cmmnt{\noindent (Bandingkan 251D.)}

\leader{254F}{Teorema} Misalkan $\langle(X_i,\Sigma_i,\mu_i)\rangle_{i\in I}$
suatu keluarga ruang probabilitas, dan misalkan $\lambda$ ukuran produk
pada $X=\prod_{i\in I}X_i$\cmmnt{ yang didefinisikan seperti dalam 254C}; misalkan
$\Lambda$ domainnya.

(a) $\lambda X=1$.

(b) Jika $E_i\in\Sigma_i$ untuk setiap $i\in I$, dan $\{i:E_i\ne X_i\}$ adalah
terhitung, maka $\prod_{i\in I}E_i\in\Lambda$, dan
$\lambda(\prod_{i\in I}E_i)=\prod_{i\in I}\mu_iE_i$. Khususnya,
$\lambda C=\theta_0C$
untuk setiap silinder terukur $C$,\cmmnt{ sebagaimana didefinisikan dalam 254A,} dan jika
$j\in I$, maka $x\mapsto x(j):X\to X_j$ merupakan fungsi pelestari ukuran melalui prapeta.

(c) $\Tensorhat_{i\in I}\Sigma_i\subseteq\Lambda$.

(d) $\lambda$ lengkap.

(e) Untuk setiap $W\in\Lambda$ dan $\epsilon>0$, terdapat suatu keluarga berhingga
$C_0,\ldots,C_n$ dari silinder terukur sedemikian sehingga
$\lambda(W\symmdiff\bigcup_{k\le n}C_k)\penalty-100\le\epsilon$.

(f) Untuk setiap $W\in\Lambda$, terdapat $W_1$,
$W_2\in\Tensorhat_{i\in I}\Sigma_i$ sedemikian sehingga
$W_1\subseteq W\subseteq W_2$ dan $\lambda(W_2\setminus W_1)=0$.

\cmmnt{\medskip

\noindent{\bf Catatan} Barangkali saya perlu berhenti sejenak untuk menafsirkan
produk $\prod_{i\in I}\mu_iE_i$. Karena semua $\mu_iE_i$ termasuk dalam
$[0,1]$, produk ini hanyalah $\inf_{J\subseteq I,J\text{
berhingga}}\prod_{i\in J}\mu_iE_i$, dengan produk kosong dianggap sama dengan $1$.
}%end of comment

\proof{ Sepanjang bukti ini, definisikan $\Cal C$, $\theta_0$
dan $\theta$ seperti dalam 254A. Saya akan menuliskan argumen yang berlaku
untuk $I$ berhingga maupun $I$ tak hingga, tetapi pada pembacaan pertama anda boleh saja
lebih suka mengandaikan bahwa $I$ tak hingga.

\medskip

{\bf (a)} Tentu saja $\lambda X=\theta X$, sehingga saya harus menunjukkan bahwa
$\theta X=1$. Karena $X$, $\emptyset\in\Cal C$ dan
$\theta_0X=\prod_{i\in I}\mu_iX_i=1$ serta $\theta_0\emptyset=0$,

\Centerline{$\theta X\le\theta_0X+\theta_0\emptyset+\ldots=1$.}

\noindent Jadi saya harus menunjukkan
bahwa $\theta X\ge 1$. \Quer\ Andaikan, jika mungkin, bahwa hal ini salah.

\medskip

\quad{\bf (i)} Terdapat suatu barisan $\sequencen{C_n}$ dalam $\Cal C$,
yang menutupi $X$, sedemikian sehingga $\sum_{n=0}^{\infty}\theta_0C_n<1$. Untuk setiap
$n\in\Bbb N$, nyatakan $C_n$ sebagai $\{x:x(i)\in E_{ni}\Forall i\in I\}$,
dengan setiap $E_{ni}\in\Sigma_i$ dan $J_n=\{i:E_{ni}\ne X_i\}$ berhingga.
Tidak satu pun $J_n$ dapat kosong, karena
$\theta_0C_n<1=\theta_0X$; tetapkan $J=\bigcup_{n\in\Bbb N}J_n$. Maka $J$
merupakan subhimpunan terhitung yang tidak kosong dari $I$. Tetapkan $K=\Bbb N$ jika $J$
tak hingga, dan $\{k:0\le k<\#(J)\}$ jika $J$ berhingga; misalkan
$k\mapsto i_k:K\to J$ suatu bijeksi.

Untuk setiap $k\in K$, tetapkan $L_k=\{i_j:j<k\}\subseteq J$, dan tetapkan
$\alpha_{nk}=\prod_{i\in I\setminus L_k}\mu_iE_{ni}$ untuk $n\in\Bbb N$,
$k\in K$. Jika $J$ berhingga, kita dapat mengidentifikasi $L_{\#(J)}$ dengan $J$,
dan menetapkan $\alpha_{n,\#(J)}=1$ untuk setiap
$n$. Kita mempunyai $\alpha_{n0}=\theta_0C_n$ untuk setiap $n$, sehingga
$\sum_{n=0}^{\infty}\alpha_{n0}<1$. Untuk $n\in\Bbb N$, $k\in K$, dan
$t\in X_{i_k}$, tetapkan

$$\eqalign{f_{nk}(t)&=\alpha_{n,k+1}\text{ jika }t\in E_{n,i_k},\cr
&=0\text{ jika tidak}.\cr}$$

\noindent Maka

\Centerline{$\int f_{nk}d\mu_{i_k}
=\alpha_{n,k+1}\mu_{i_k}E_{n,i_k}=\alpha_{nk}$.}

\medskip

\quad{\bf (ii)} Pilih $t_k\in X_{i_k}$ secara induktif, untuk $k\in K$, sebagai
berikut. Hipotesis induksinya adalah
$\sum_{n\in M_k}\alpha_{nk}<1$, dengan
$M_k=\{n:n\in \Bbb N$, $t_j\in E_{n,i_j}\Forall j<k\}$; tentu saja
$M_0=\Bbb N$, sehingga induksi dapat dimulai. Karena

\Centerline{$1>\sum_{n\in M_k}\alpha_{nk}
=\sum_{n\in M_k}\int f_{nk}d\mu_{i_k}
=\int(\sum_{n\in M_k}f_{nk})d\mu_{i_k}$}

\noindent (menurut teorema B. Levi), harus terdapat $t_k\in X_{i_k}$ sedemikian
sehingga $\sum_{n\in M_k}f_{nk}(t_{k})<1$. Untuk pilihan
$t_k$ semacam itu, $\alpha_{n,k+1}=f_{nk}(t_k)$ untuk setiap $n\in M_{k+1}$, sehingga
$\sum_{n\in M_{k+1}}\alpha_{n,k+1}<1$, dan induksi berlanjut,
kecuali jika $J$ berhingga dan $k+1=\#(J)$. Dalam kasus terakhir ini kita harus
mempunyai $M_{\#(J)}=\emptyset$, karena $\alpha_{n,\#(J)}=1$ untuk setiap $n$.

\medskip

\quad{\bf (iii)} Jika $J$ tak hingga, kita memperoleh suatu barisan lengkap
$\sequence{k}{t_k}$; jika $J$ berhingga, kita hanya memperoleh barisan berhingga
$\langle t_k\rangle_{k<\#(J)}$. Dalam kedua kasus, terdapat $x\in X$
sedemikian sehingga $x(i_k)=t_k$ untuk setiap $k\in K$. Sekarang harus terdapat
$m\in\Bbb N$ sedemikian sehingga $x\in C_m$. Karena $J_m=\{i:E_{mi}\ne X_i\}$
berhingga, terdapat $k\in\Bbb N$ sedemikian sehingga $J_m\subseteq L_k$
(dengan mengizinkan $k=\#(J)$ jika $J$ berhingga). Sekarang $m\in M_k$, sehingga
sesungguhnya tidak mungkin $k=\#(J)$, dan
$\alpha_{mk}=1$, sehingga $\sum_{n\in M_k}\alpha_{nk}\ge 1$, bertentangan dengan
hipotesis induksi.\ \Bang

Kontradiksi ini menunjukkan bahwa $\theta X=1$.

\medskip

{\bf (b)(i)} Saya membahas kasus khusus terlebih dahulu.
Andaikan $j\in I$ dan
$E\in\Sigma_j$, serta misalkan $C\in\Cal C$; tetapkan $W=\{x:x\in X$, $x(j)\in
E\}$;
maka $C\cap W$ dan $C\setminus W$ keduanya termasuk dalam $\Cal C$, dan $\theta_0
C=\theta_0(C\cap W)+\theta_0(C\setminus W)$. \Prf\ Jika
$C=\prod_{i\in I}C_i$, dengan $C_i\in\Sigma_i$ untuk setiap $i$, maka
$C\cap W=\prod_{i\in I}C_i'$, dengan $C_i'=C_i$ jika $i\ne j$, dan
$C_j'=C_j\cap E$; demikian pula,
$C\setminus W=\prod_{i\in I}C_i''$, dengan $C_i''=C_i$ jika $i\ne j$, dan
$C_j''=C_j\setminus E$. Jadi keduanya termasuk dalam $\Cal C$, dan

\Centerline{$\theta_0(C\cap W)+\theta_0(C\setminus W)
=(\mu_j(C_j\cap E)+\mu_j(C_j\setminus E))\prod_{i\ne j}\mu_iC_i
=\prod_{i\in I}\mu_iC_i
=\theta_0C$.  \Qed}

\medskip

\quad{\bf (ii)} Sekarang andaikan $A\subseteq X$ sebarang himpunan, dan
$\epsilon>0$. Maka terdapat suatu barisan $\sequencen{C_n}$ dalam $\Cal C$
sedemikian sehingga $A\subseteq\bigcup_{n\in\Bbb N}C_n$ dan
$\sum_{n=0}^{\infty}\theta_0C_n\le\theta A+\epsilon$. Dalam hal ini

\Centerline{$A\cap W\subseteq\bigcup_{n\in\Bbb N}C_n\cap W$,\quad
$A\setminus W\subseteq\bigcup_{n\in\Bbb N}C_n\setminus W,$}

\noindent sehingga

\Centerline{$\theta(A\cap W)\le\sum_{n=0}^{\infty}\theta_0(C_n\cap
W),\quad
\theta(A\setminus W)\le\sum_{n=0}^{\infty}\theta_0(C_n\setminus W),$}

\noindent dan

\Centerline{$\theta(A\cap W)+\theta(A\setminus W)
\le\sum_{n=0}^{\infty}\theta_0(C_n\cap W)+\theta_0(C_n\setminus W)
=\sum_{n=0}^{\infty}\theta_0C_n
\le\theta A+\epsilon.$}
\noindent Karena $\epsilon$ sembarang, $\theta(A\cap W)+\theta(A\setminus
W)\le\theta A$; karena $A$ sembarang, $W\in\Lambda$.

\medskip

\quad{\bf (iii)}
Berikutnya saya menunjukkan bahwa jika $J\subseteq I$ berhingga dan $C_i\in\Sigma_i$ untuk
setiap $i\in J$, serta $C=\{x:x\in X$, $x(i)\in C_i\Forall i\in J\}$,
maka $C\in\Lambda$ dan $\lambda C=\prod_{i\in J}\mu_iC_i$. \Prf\
Lakukan induksi terhadap $\#(J)$. Jika $\#(J)=0$, yaitu $J=\emptyset$, maka $C=X$
dan ini adalah bagian (a). Untuk langkah induksi menuju $\#(J)=n+1$, ambil sebarang
$j\in J$ dan tetapkan $J'=J\setminus\{j\}$,

\Centerline{$C'=\{x:x\in X$, $x(i)\in C_i\Forall i\in J'\},$}

\Centerline{$C''=C'\setminus C=\{x:x\in C'$, $x(j)\in X_j\setminus
C_j\}.$}

\noindent Maka $C$, $C'$, dan $C''$ semuanya termasuk dalam $\Cal C$, dan $\theta_0
C'=\prod_{i\in J'}\mu_iC_i=\alpha$, katakanlah, $\theta_0C=\alpha\mu_jC_j$,
$\theta_0C''=\alpha(1-\mu_jC_j)$. Selain itu, menurut hipotesis
induksi, $C'\in\Lambda$ dan $\alpha=\lambda C'=\theta C'$. Jadi
$C=C'\cap\{x:x(j)\in C_j\}\in\Lambda$ menurut (ii), dan $C''=C'\setminus
C\in\Lambda$.

Kita tentu mempunyai $\lambda C=\theta C\le\theta_0C$, $\lambda
C''\le\theta_0C''$; tetapi juga

\Centerline{$\alpha=\lambda C'=\lambda C+\lambda C''\le\theta_0
C+\theta_0C''=\alpha,$}

\noindent sehingga sesungguhnya

\Centerline{$\lambda C=\theta_0 C=\alpha\mu_jC_j
=\prod_{i\in J}\mu_iC_i$,}

\noindent dan induksi berlanjut.\ \Qed

\medskip

\quad{\bf (iv)} Sekarang mari kembali ke kasus umum suatu himpunan $W$ yang
berbentuk $\prod_{i\in I}E_i$, dengan $E_i\in\Sigma_i$ untuk setiap $i$, dan
$K=\{i:E_i\ne X_i\}$ terhitung. Jika $K$ berhingga, maka
$W=\{x:x(i)\in E_i\Forall i\in K\}$, sehingga $W\in\Lambda$ dan

\Centerline{$\lambda W
=\prod_{i\in K}\mu_iE_i=\prod_{i\in I}\mu_iE_i$.}

\noindent Jika tidak, misalkan
$\sequencen{i_n}$ suatu enumerasi dari $K$. Untuk setiap $n\in\Bbb N$, tetapkan
$W_n=\{x:x\in X$, $x(i_k)\in E_{i_k}\Forall k\le n\}$; kita telah mengetahui
bahwa $W_n\in\Lambda$ dan bahwa
$\lambda W_n=\prod_{k=0}^n\mu_{i_k}E_{i_k}$. Namun $\sequencen{W_n}$ merupakan
barisan tak menaik dengan irisan $W$, sehingga $W\in\Lambda$ dan

\Centerline{$\lambda W=\lim_{n\to\infty}\lambda W_n
=\prod_{i\in K}\mu_iE_i
=\prod_{i\in I}\mu_iE_i.$}

\medskip

{\bf (c)} merupakan akibat langsung dari (b) dan definisi
$\Tensorhat_{i\in I}\Sigma_i$.

\medskip

{\bf (d)} Karena $\lambda$ dibangun dengan metode \Caratheodory, ukuran ini
harus lengkap.

\medskip

{\bf (e)} Misalkan $\sequencen{C_n}$ suatu barisan dalam $\Cal C$ sedemikian sehingga
$W\subseteq\bigcup_{n\in\Bbb N}C_n$ dan
$\sum_{n=0}^{\infty}\theta_0C_n\le\theta W+\bover12\epsilon$. Tetapkan
$V=\bigcup_{n\in\Bbb N}C_n$; menurut (b), $V\in\Lambda$. Misalkan $n\in\Bbb N$
sedemikian
sehingga $\sum_{i=n+1}^{\infty}\theta_0C_i\le\bover12\epsilon$, dan
tinjau $W'=\bigcup_{k\le n}C_k$.
Karena $V\setminus W'\subseteq\bigcup_{i>n}C_i$,

$$\eqalign{\lambda(W\symmdiff W')
&\le\lambda(V\setminus W)+\lambda(V\setminus W')
=\lambda V-\lambda W+\lambda(V\setminus W')
=\theta V-\theta W+\theta(V\setminus W')\cr
&\le\sum_{i=0}^{\infty}\theta_0C_i-\theta W+\sum_{i=n+1}\theta_0C_i
\le\Bover12\epsilon+\Bover12\epsilon
=\epsilon.\cr}$$

\medskip

{\bf (f)(i)} Jika $W\in\Lambda$ dan $\epsilon>0$, terdapat suatu
$V\in\Tensorhat_{i\in I}\Sigma_i$ sedemikian sehingga $W\subseteq V$ dan $\lambda
V\le\lambda W+\epsilon$. \Prf\ Misalkan $\sequencen{C_n}$ suatu barisan dalam
$\Cal C$ sedemikian sehingga $W\subseteq\bigcup_{n\in\Bbb N}C_n$ dan
$\sum_{n=0}^{\infty}\theta_0C_n\le\theta W+\epsilon$. Maka
$C_n\in\Tensorhat_{i\in I}\Sigma_i$ untuk setiap $n$, sehingga
$V=\bigcup_{n\in\Bbb N}C_n\in\Tensorhat_{i\in I}\Sigma_i$. Sekarang
$W\subseteq V$, dan

\Centerline{$\lambda V=\theta V\le\sum_{n=0}^{\infty}\theta_0C_n
\le\theta W+\epsilon=\lambda W+\epsilon$.\ \Qed}

\medskip

\quad{\bf (ii)} Sekarang, untuk $W\in\Lambda$, misalkan $\sequencen{V_n}$ suatu
barisan himpunan dalam $\Tensorhat_{i\in I}\Sigma_i$ sedemikian sehingga $W\subseteq
V_n$ dan $\lambda V_n\le\lambda W+2^{-n}$ untuk setiap $n$; maka
$W_2=\bigcap_{n\in\Bbb N}V_n$ termasuk dalam $\Tensorhat_{i\in I}\Sigma_i$
dan
$W\subseteq W_2$ dan $\lambda W_2=\lambda W$. Dengan cara yang sama, terdapat
$W'_2\in\Tensorhat_{i\in I}\Sigma_i$ sedemikian sehingga $X\setminus W\subseteq
W'_2$ dan $\lambda W'_2=\lambda(X\setminus W)$, sehingga kita dapat mengambil
$W_1=X\setminus W'_2$ untuk menyelesaikan bukti.
}

\leader{254G}{}\cmmnt{ Berikut ini adalah sifat mendasar, bahkan
sifat penentu, dari ukuran produk. (Bandingkan 251L.)

\medskip

\noindent}{\bf Teorema} Misalkan $\familyiI{(X_i,\Sigma_i,\mu_i)}$ suatu
keluarga
ruang probabilitas dengan produk $(X,\Lambda,\lambda)$. Misalkan
$(Y,\Tau,\nu)$ suatu ruang probabilitas lengkap dan $\phi:Y\to X$ suatu
fungsi.
Andaikan bahwa $\nu^*\phi^{-1}[C]\le\lambda C$ untuk setiap silinder terukur
$C\subseteq X$. Maka $\phi$ bersifat \imp.
Khususnya, $\phi$ bersifat \imp\ jika dan hanya jika $\phi^{-1}[C]\in\Tau$ dan
$\nu\phi^{-1}[C]=\lambda C$ untuk setiap silinder
terukur $C\subseteq X$.

\cmmnt{\medskip

\noindent{\bf Catatan} Dengan $\nu^*$ saya maksudkan ukuran luar yang lazim, yang
didefinisikan dari $\nu$ seperti dalam \S132.
}

\proof{{\bf (a)} Pertama-tama perhatikan bahwa, dengan menuliskan $\theta$ untuk ukuran luar
dalam 254A, $\nu^*\phi^{-1}[A]\le\theta A$ untuk setiap $A\subseteq X$.
\Prf\ Untuk $\epsilon>0$, terdapat suatu barisan $\sequencen{C_n}$ dari
silinder terukur sedemikian sehingga $A\subseteq\bigcup_{n\in\Bbb N}C_n$ dan
$\sum_{n=0}^{\infty}\theta_0C_n\le\theta A+\epsilon$, dengan $\theta_0$
fungsional dalam 254A. Namun kita mengetahui bahwa $\theta_0C=\lambda C$ untuk
setiap silinder terukur $C$ (254Fb), sehingga

\Centerline{$\nu^*\phi^{-1}[A]
\le\nu^*(\bigcup_{n\in\Bbb N}\phi^{-1}[C_n])
\le\sum_{n=0}^{\infty}\nu^*\phi^{-1}[C_n]
\le\sum_{n=0}^{\infty}\lambda C_n
\le\theta A+\epsilon$.}

\noindent Karena $\epsilon$ sembarang, $\nu^*\phi^{-1}[A]\le\theta A$.\
\Qed

\medskip

{\bf (b)} Sekarang ambil sebarang $W\in\Lambda$. Maka terdapat $F$, $F'\in\Tau$
sedemikian sehingga

\Centerline{$\phi^{-1}[W]\subseteq F$, \quad$\phi^{-1}[X\setminus
W]\subseteq F'$,}

\Centerline{$\nu F=\nu^*\phi^{-1}[W]\le\theta W=\lambda W$,
\quad$\nu F'\le\lambda[X\setminus W]$.}

\noindent Kita mempunyai

\Centerline{$F\cup F'\supseteq\phi^{-1}[W]\cup\phi^{-1}[X\setminus
W]=Y$,}

\noindent sehingga

\Centerline{$\nu(F\cap F')=\nu F+\nu F'-\nu(F\cup F')
\le\lambda W+\lambda(X\setminus W)-1=0$.}

\noindent Sekarang

\Centerline{$F\setminus\phi^{-1}[W]\subseteq F\cap\phi^{-1}[X\setminus
W]\subseteq F\cap F'$}

\noindent terabaikan terhadap $\nu$. Karena $\nu$ lengkap,
$F\setminus\phi^{-1}[W]\in\Tau$ dan
$\phi^{-1}[W]=F\setminus(F\setminus\phi^{-1}[W])$ termasuk dalam $\Tau$.
Selain itu,

\Centerline{$1=\nu F+\nu F'\le\lambda W+\lambda(X\setminus W)=1$,}

\noindent sehingga harus berlaku $\nu F=\lambda W$; tetapi ini berarti
bahwa $\nu\phi^{-1}[W]=\lambda W$. Karena $W$ sembarang, $\phi$ bersifat \imp.
}%end of proof of 254G

\leader{254H}{Korolari} Misalkan
$\langle(X_i,\Sigma_i,\mu_i)\rangle_{i\in I}$
dan $\langle(Y_i,\Tau_i,\nu_i)\rangle_{i\in I}$ dua keluarga
ruang probabilitas, dengan produk $(X,\Lambda,\lambda)$ dan
$(Y,\Lambda',\lambda')$. Andaikan bahwa untuk setiap $i\in I$ diberikan
suatu fungsi \imp\ $\phi_i:X_i\to Y_i$. Tetapkan
$\phi(x)=\familyiI{\phi_i(x(i))}$ untuk $x\in X$. Maka $\phi:X\to Y$ bersifat
\imp.

\proof{ Jika $C=\prod_{i\in I}C_i$ merupakan silinder terukur dalam $Y$, maka
$\phi^{-1}[C]=\prod_{i\in I}\phi_i^{-1}[C_i]$ merupakan silinder terukur
dalam
$X$, dan

\Centerline{$\lambda\phi^{-1}[C]
=\prod_{i\in I}\mu_i\phi_i^{-1}[C_i]
=\prod_{i\in I}\nu_i C_i
=\lambda'C$.}

\noindent Karena $\lambda$ merupakan ukuran probabilitas lengkap, 254G memberi
tahu kita bahwa $\phi$ bersifat \imp.
}%end of proof of 254H

\leader{254I}{}\cmmnt{ Sebagai padanan 251T, kita mempunyai hasil berikut.

\medskip

\noindent}{\bf Proposisi} Misalkan
$\langle(X_i,\Sigma_i,\mu_i)\rangle_{i\in I}$ suatu keluarga ruang probabilitas,
$\lambda$ ukuran produk
pada $X=\prod_{i\in I}X_i$, dan $\Lambda$ domainnya. Maka $\lambda$ juga
merupakan produk dari
pelengkapan $\hat\mu_i$ masing-masing $\mu_i$\cmmnt{ (212C)}.

\proof{ Tuliskan $\hat\lambda$ untuk produk dari $\hat\mu_i$, dan
$\hat\Lambda$ untuk domainnya. (i) Peta identitas dari $X_i$ ke
dirinya sendiri bersifat \imp\ jika dipandang sebagai peta dari $(X_i,\hat\mu_i)$ ke
$(X_i,\mu_i)$, sehingga
peta identitas pada $X$ bersifat \imp\ jika dipandang sebagai peta dari
$(X,\hat\lambda)$ ke $(X,\lambda)$, menurut 254H; yaitu,
$\Lambda\subseteq\hat\Lambda$ dan $\lambda=\hat\lambda\restr\Lambda$.
(ii) Jika $C$ merupakan silinder terukur untuk $\langle\hat\mu_i\rangle_{i\in
I}$, yaitu, $C=\prod_{i\in I}C_i$ dengan $C_i\in\hat\Sigma_i$ untuk setiap
$i$ dan $\{i:C_i\ne X_i\}$ adalah
berhingga, maka untuk setiap $i\in I$ kita dapat menemukan suatu $C'_i\in\Sigma_i$ sedemikian sehingga
$C_i\subseteq C'_i$ dan $\mu_iC'_i=\hat\mu_iC_i$; dengan menetapkan
$C'=\prod_{i\in I}C'_i$, kita memperoleh

\Centerline{$\lambda^*C\le\lambda C'=\prod_{i\in I}\mu_iC'_i
=\prod_{i\in I}\hat\mu_iC_i=\hat\lambda C$.}

\noindent Menurut 254G, $\lambda W$ harus terdefinisi dan sama dengan
$\hat\lambda W$ kapan pun $W\in\hat\Lambda$. Dengan menggabungkan hal ini
dengan (i), kita melihat bahwa $\lambda=\hat\lambda$.
}%end of proof of 254I

\leader{254J}{}{\bf Ukuran produk pada $\{0,1\}^I$ (a)} Barangkali
contoh terpenting di antara semua ukuran produk tak hingga adalah kasus
ketika setiap faktor $X_i$ hanyalah $\{0,1\}$ dan setiap $\mu_i$ merupakan
ukuran probabilitas ‘koin adil’, yang ditetapkan oleh

\Centerline{$\mu_i\{0\}=\mu_i\{1\}=\Bover12$.}

\noindent Dalam kasus ini, produk $X=\{0,1\}^I$ mempunyai suatu keluarga
$\langle E_i\rangle_{i\in I}$ dari himpunan terukur sedemikian sehingga,
dengan menuliskan $\lambda$ untuk ukuran produk pada $X$,

\Centerline{$\lambda(\bigcap_{i\in J}E_i)=2^{-\#(J)}$ jika $J\subseteq I$
berhingga.}

\noindent\prooflet{(Cukup ambil $E_i=\{x:x(i)=1\}$ untuk setiap $i$.) }Saya
akan menyebut
$\lambda$ ini {\bf ukuran biasa} pada $\{0,1\}^I$. Perhatikan bahwa jika $I$
berhingga, maka $\lambda\{x\}=2^{-\#(I)}$ untuk setiap
$x\in X$\cmmnt{ (dengan menggunakan
254Fb)}. Di sisi lain, jika $I$ tak hingga, maka $\lambda\{x\}=0$
untuk setiap $x\in X$\prooflet{ (karena, sekali lagi dengan menggunakan 254Fb,
$\lambda^*\{x\}\le 2^{-n}$ untuk setiap $n$)}.

\spheader 254Jb Terdapat bijeksi alami antara $\{0,1\}^I$ dan
$\Cal PI$, yang memasangkan $x\in\{0,1\}^I$ dengan $\{i:i\in I$, $x(i)=1\}$.
Jadi kita memperoleh suatu ukuran baku $\tilde\lambda$ pada $\Cal PI$, yang akan
saya sebut {\bf ukuran biasa pada $\Cal PI$}. Perhatikan bahwa untuk sebarang
$b\subseteq I$ berhingga dan sebarang $c\subseteq b$, kita mempunyai

\Centerline{$\tilde\lambda\{a:a\cap b=c\}=\lambda\{x:x(i)=1$ untuk $i\in
c$, $x(i)=0$ untuk $i\in b\setminus c\}=2^{-\#(b)}$.}

\spheader 254Jc Tentu saja kita dapat menerapkan 254G pada ukuran-ukuran ini; jika
$(Y,\Tau,\nu)$ merupakan ruang probabilitas lengkap, suatu fungsi
$\phi:Y\to\{0,1\}^I$ bersifat \imp\ jika dan hanya jika

\Centerline{$\nu\{y:y\in Y$, $\phi(y)\restr J=z\}=2^{-\#(J)}$}

\noindent kapan pun $J\subseteq I$ berhingga dan
$z\in\{0,1\}^J$\prooflet{; ini karena silinder terukur dalam
$\{0,1\}^I$ tepat merupakan himpunan berbentuk $\{x:x\restr J=z\}$ dengan
$J\subseteq I$ berhingga}.

\spheader 254Jd\dvAformerly{2{}54Xe}
Definisikan penjumlahan pada $X$ dengan menetapkan
$(x+y)(i)=x(i)+_2y(i)$ untuk setiap $i\in I$, $x$, $y\in X$, dengan
$0+_20=1+_21=0$, $0+_21=1+_20=1$. Jika $y\in X$, peta
$x\mapsto x+y:X\to X$ bersifat\cmmnt{ \imp. \prooflet{\Prf\ Jika $J\subseteq I$
berhingga dan $z\in\{0,1\}^J$, tetapkan $z'=\family{j}{J}{z(j)+_2y(j)}$; maka

\Centerline{$\lambda\{x:(x+y)\restr J=z\}
=\lambda\{x:x\restr J=z'\}=2^{-\#(J)}$.}

\noindent Karena $J$ sembarang, (c) memberi tahu kita bahwa $x\mapsto x+y$ bersifat \imp.\
\Qed} Sekarang, karena

\Centerline{$(x+y)+y=x+(y+y)=x+0=x$}

\noindent untuk setiap $x$, peta $x\mapsto x+y:X\to X$ bersifat bijektif dan
sama dengan inversnya, sehingga peta itu sebenarnya} suatu automorfisme ruang ukur
dari $(X,\lambda)$.

%\def\spheader#1#2#3#4#5{\header{#1#2#3#4#5}{\bf (#5)}}
%\def\header#1{
%     \wheader{#1}{10}{4}{4}{24pt}}

\header{254Je}{\bf *(e)}\dvAnew{2016}
Karena semua faktor $(X_i,\mu_i)$ sama, kita
mempunyai kelas lain automorfisme dari $(X,\lambda)$, yang berkaitan dengan
permutasi $I$. Jika $\pi:I\to I$ sebarang permutasi, maka
kita mempunyai fungsi yang berkaitan $x\mapsto x\pi:X\to X$. \cmmnt{Jika
$J\subseteq I$ berhingga dan $z\in\{0,1\}^J$, tetapkan $J'=\pi[J]$ dan
$z'=z\pi^{-1}\in\{0,1\}^{J'}$; maka

\Centerline{$\lambda\{x:(x\pi)\restr J=z\}=\lambda\{x:x\restr J'=z'\}
=2^{-\#(J')}=2^{-\#(J)}$.}

\noindent Jadi $x\mapsto x\pi$ bersifat \imp. Kali ini, inversnya adalah
$x\mapsto x\pi^{-1}$, yang juga bersifat \imp; sehingga} $x\mapsto x\pi$
merupakan automorfisme ruang ukur.

\leader{254K}{}\cmmnt{Dalam kasus $I$ tak hingga terhitung, kita
mempunyai hubungan yang sangat penting
antara ukuran produk biasa pada $\{0,1\}^I$ dan
ukuran Lebesgue pada $[0,1]$.

\medskip

\noindent}{\bf Proposisi} Misalkan $\lambda$ ukuran biasa pada
$X=\{0,1\}^{\Bbb N}$, dan misalkan $\mu$ ukuran Lebesgue pada $[0,1]$;
tuliskan $\Lambda$ untuk domain $\lambda$ dan $\Sigma$ untuk domain
$\mu$.

(i) Untuk $x\in X$, tetapkan $\phi(x)=\sum_{i=0}^{\infty}2^{-i-1}x(i)$. Maka

\qquad $\phi^{-1}[E]\in\Lambda$ dan $\lambda\phi^{-1}[E]=\mu E$
untuk setiap $E\in\Sigma$;

\qquad $\phi[F]\in\Sigma$ dan $\mu\phi[F]=\lambda F$ untuk
setiap $F\in\Lambda$.

(ii) Terdapat suatu bijeksi $\tilde\phi:X\to[0,1]$ yang sama dengan
$\phi$ kecuali pada terhitung banyak titik, dan sebarang bijeksi semacam itu merupakan
isomorfisme antara $(X,\Lambda,\lambda)$ dan $([0,1],\Sigma,\mu)$.

\proof{{\bf (a)} Hal pertama yang perlu diperhatikan ialah bahwa $\phi$ sendiri
hampir merupakan bijeksi. Dengan menetapkan

\Centerline{$H=\{x:x\in X$, $\exists\,m\in\Bbb
N$, $x(i)=x(m)\Forall i\ge m\}$,}

\Centerline{$H'=\{2^{-n}k:n\in\Bbb N$, $k\le 2^n\}$,}

\noindent maka $H$ dan $H'$ terhitung dan $\phi\restr X\setminus H$
merupakan bijeksi antara $X\setminus H$ dan $[0,1]\setminus H'$. (Untuk
$t\in[0,1]\setminus H'$, $\phi^{-1}(t)$ adalah ekspansi biner dari $t$.)
Karena $H$ dan $H'$ tak hingga terhitung, terdapat bijeksi
di antara keduanya; dengan menggabungkan bijeksi ini dengan $\phi\restr X\setminus H$, kita mempunyai
bijeksi antara $X$ dan $[0,1]$ yang sama dengan $\phi$ kecuali pada terhitung
banyak titik. Untuk sisa bukti ini, misalkan $\tilde\phi$ sebarang bijeksi semacam itu.
Misalkan $M$ himpunan terhitung $\{x:x\in
X$, $\phi(x)\ne\tilde\phi(x)\}$, dan $N$ himpunan terhitung
$\phi[M]\cup\tilde\phi[M]$; maka
$\phi[A]\symmdiff\tilde\phi[A]\subseteq N$ untuk setiap $A\subseteq X$.

\medskip

{\bf (b)} Untuk melihat bahwa $\lambda\tilde\phi^{-1}[E]$ ada dan sama dengan
$\mu E$ untuk setiap $E\in\Sigma$, saya meninjau berturut-turut himpunan
$E$ yang makin rumit.

\medskip

\quad\grheada\ Jika $E=\{t\}$, maka
$\lambda\tilde\phi^{-1}[E]=\lambda\{\tilde\phi^{-1}(t)\}$ ada dan
bernilai nol.

\medskip

\quad\grheadb\ Jika $E$ berbentuk
$\coint{2^{-n}k,2^{-n}(k+1)}$, dengan $n\in\Bbb N$ dan $0\le k<2^n$, maka
$\phi^{-1}[E]$ berbeda paling banyak dua titik dari suatu himpunan berbentuk
$\{x:x(i)=z(i)\Forall i<n\}$, sehingga $\tilde\phi^{-1}[E]$ berselisih
dengan himpunan ini hanya pada suatu himpunan terhitung, dan

\Centerline{$\lambda\tilde\phi^{-1}[E]=2^{-n}=\mu E$.}

\medskip

\quad\grheadc\ Jika $E$ berbentuk
$\coint{2^{-n}k,2^{-n}l}$, dengan $n\in\Bbb N$ dan $0\le k<l\le 2^n$,
maka
\discrcenter{468pt}
{$E=\bigcup_{k\le i<l}\coint{2^{-n}i,2^{-n}(i+1)}$, }sehingga

\Centerline{$\lambda\tilde\phi^{-1}[E]=2^{-n}(l-k)=\mu E$.}

\medskip

\quad\grheadd\ Jika $E$ berbentuk $\coint{t,u}$, dengan
$0\le t<u\le 1$, maka untuk setiap $n\in\Bbb N$, tetapkan
$k_n=\lfloor 2^nt\rfloor$, yaitu bagian bulat
dari $2^nt$, $l_n=\lfloor 2^nu\rfloor$, dan
$E_n=\coint{2^{-n}(k_n+1),2^{-n}l_n}$; maka $\sequencen{E_n}$ merupakan
barisan tak menurun dan $\bigcup_{n\in\Bbb N}E_n$ adalah $\ooint{t,u}$.
Jadi (dengan menggunakan ($\alpha$))

$$\eqalign{\lambda\tilde\phi^{-1}[E]
&=\lambda\tilde\phi^{-1}[\bigcup_{n\in\Bbb N}E_n]
=\lim_{n\to\infty}\lambda\tilde\phi^{-1}[E_n]\cr
&=\lim_{n\to\infty}\mu E_n
=\mu E.\cr}$$

\medskip

\quad\grheade\ Jika $E\in\Sigma$, maka untuk sebarang
$\epsilon>0$, terdapat suatu barisan $\sequencen{I_n}$ dari
subinterval setengah terbuka dari $\coint{0,1}$ sedemikian sehingga
$E\setminus\{1\}\subseteq\bigcup_{n\in\Bbb N}I_n$ dan
$\sum_{n=0}^{\infty}\mu I_n\le\mu E+\epsilon$; sekarang
$\tilde\phi^{-1}[E]
\subseteq \{\tilde\phi^{-1}(1)\}\cup\bigcup_{n\in\Bbb N}\tilde\phi^{-1}[I_n]$,
sehingga

\Centerline{$\lambda^*\tilde\phi^{-1}[E]
\le\lambda(\bigcup_{n\in\Bbb N}\tilde\phi^{-1}[I_n])
\le \sum_{n=0}^{\infty}\lambda\tilde\phi^{-1}[I_n]
=\sum_{n=0}^{\infty}\mu I_n
\le\mu E+\epsilon$.}

\noindent Karena $\epsilon$ sembarang,
$\lambda^*\tilde\phi^{-1}[E]\le\mu E$, dan
terdapat $V\in\Lambda$ sedemikian sehingga $\tilde\phi^{-1}[E]\subseteq V$ dan
$\lambda V\le\mu E$.

\medskip

\quad\grheadz\ Dengan cara yang sama, terdapat $V'\in\Lambda$ sedemikian
sehingga $V'\supseteq\tilde\phi^{-1}[[0,1]\setminus E]$ dan
$\lambda V'\le\mu([0,1]\setminus E)$. Sekarang $V\cup V'=X$, sehingga

\Centerline{$\lambda(V\cap V')=\lambda V+\lambda V'-\lambda(V\cup V')
\le\mu E+(1-\mu E)-1=0$}

\noindent dan

\Centerline{$\tilde\phi^{-1}[E]=(X\setminus V')
\cup(V\cap V'\cap\tilde\phi^{-1}[E])$}

\noindent termasuk dalam $\Lambda$, dengan

\Centerline{$\lambda\tilde\phi^{-1}[E]\le\lambda V\le\mu E$;}

\noindent pada saat yang sama,

\Centerline{$1-\lambda\tilde\phi^{-1}[E]\le\lambda V'\le 1-\mu E$}

\noindent sehingga $\lambda\tilde\phi^{-1}[E]=\mu E$.

\medskip

{\bf (c)} Sekarang andaikan $C\subseteq X$ suatu silinder terukur dengan bentuk khusus
$\{x:x(0)=\epsilon_0,\ldots,x(n)=\epsilon_n\}$ untuk suatu
$\epsilon_0,\ldots,\epsilon_n\in\{0,1\}$. Maka
$\phi[C]=[t,t+2^{-n-1}]$ dengan $t=\sum_{i=0}^n2^{-i-1}\epsilon_i$, sehingga
$\mu\phi[C]=\lambda C$. Karena
$\tilde\phi[C]\symmdiff\phi[C]\subseteq N$ terhitung,
$\mu\tilde\phi[C]=\lambda C$.

Jika $C\subseteq X$ sebarang silinder terukur, maka silinder ini berbentuk
$\{x:x\restr J=z\}$ untuk suatu $J\subseteq\Bbb N$ berhingga; dengan mengambil $n$ yang
cukup besar sehingga $J\subseteq\{0,\ldots,n\}$, $C$ dapat dinyatakan sebagai gabungan
saling lepas dari $2^{n+1-\#(J)}$ himpunan berbentuk seperti yang baru saja dibahas, yaitu
himpunan-himpunan dengan $\epsilon_i=z(i)$ untuk $i\in J$. Dengan menjumlahkan ukurannya,
kita kembali memperoleh $\mu\tilde\phi[C]=\lambda C$. Sekarang 254G memberi tahu kita bahwa
$\tilde\phi^{-1}:[0,1]\to X$ bersifat \imp, yaitu, $\tilde\phi[W]$
terukur Lebesgue, dengan ukuran $\lambda W$, untuk setiap $W\in\Lambda$.

Dengan menggabungkan hal ini dengan (b), $\tilde\phi$ harus merupakan isomorfisme
antara $(X,\Lambda,\lambda)$ dan $([0,1],\Sigma,\mu)$, sebagaimana ditegaskan dalam
(ii) dari proposisi.

\medskip

{\bf (d)} Mengenai (i), jika $E\in\Sigma$, maka
$\phi^{-1}[E]\symmdiff\tilde\phi^{-1}[E]\subseteq M$ terhitung, sehingga
$\lambda\phi^{-1}[E]=\lambda\tilde\phi^{-1}[E]=\mu E$. Sedangkan jika
$W\in\Lambda$, $\phi[W]\symmdiff\tilde\phi[W]\subseteq N$ terhitung,
sehingga $\mu\phi[W]=\mu\tilde\phi[W]=\lambda W$.
}%end of proof of 254K

\leader{254L}{\dvrocolon{Subruang}}\cmmnt{ Sebagaimana dalam 251Q, kita dapat
meninjau produk dari
ukuran-ukuran subruang.   Bentuk hasilnya lebih sederhana
karena dalam konteks sekarang kita terbatas pada
ukuran-ukuran probabilitas.

\medskip

\noindent}{\bf Teorema} Misalkan
$\langle(X_i,\Sigma_i,\mu_i)\rangle_{i\in I}$
suatu keluarga ruang probabilitas, dan $(X,\Lambda,\lambda)$ produk
mereka.

(a) Untuk setiap $i\in I$, misalkan $A_i\subseteq X_i$ suatu himpunan dengan
ukuran luar penuh, dan tuliskan $\tilde\mu_i$ untuk ukuran subruang pada
$A_i$\cmmnt{ (214B)}.
Misalkan $\tilde\lambda$ ukuran produk pada $A=\prod_{i\in I}A_i$.
Maka $\tilde\lambda$ adalah ukuran subruang pada $A$ yang diinduksi oleh
$\lambda$.

(b) $\lambda^*(\prod_{i\in I}A_i)=\prod_{i\in I}\mu_i^*A_i$ kapan pun
$A_i\subseteq X_i$ untuk setiap $i$.

\proof{{\bf (a)} Tuliskan $\lambda_A$ untuk ukuran subruang pada $A$
yang didefinisikan dari $\lambda$, dan $\Lambda_A$ untuk domainnya;  tuliskan
$\tilde\Lambda$ untuk domain $\tilde\lambda$.

\medskip

\quad{\bf (i)} Misalkan $\phi:A\to X$ peta identitas.   Jika $C\subseteq
X$ suatu silinder terukur, katakanlah $C=\prod_{i\in I}C_i$ dengan
$C_i\in\Sigma_i$ untuk setiap $i$, maka $\phi^{-1}[C]=\prod_{i\in I}(C_i\cap
A_i)$ merupakan silinder terukur dalam $A$, dan

\Centerline{$\tilde\lambda\phi^{-1}[C]=\prod_{i\in I}\tilde\mu_i(C_i\cap
A_i)\le\prod_{i\in I}\mu_iC_i=\lambda C$.}

\noindent Menurut 254G, $\phi$ bersifat \imp, yaitu $\tilde\lambda(A\cap
W)=\lambda W$ untuk setiap $W\in\Lambda$.   Tetapi ini
berarti bahwa $\tilde\lambda V$ terdefinisi dan sama dengan
$\lambda_AV=\lambda^*V$ untuk setiap $V\in\Lambda_A$, karena
untuk setiap $V$ demikian terdapat $W\in\Lambda$ sedemikian sehingga $V=A\cap
W$ dan $\lambda W=\lambda_AV$.   Khususnya, $\lambda_AA=1$.

\medskip

\quad{\bf (ii)} Sekarang pandang $\phi$ sebagai fungsi dari ruang ukur
$(A,\Lambda_A,\lambda_A)$ ke $(A,\tilde\Lambda,\tilde\lambda)$.   Jika $D$
merupakan silinder terukur
dalam $A$, kita dapat menyatakannya sebagai $\prod_{i\in I}D_i$ dengan setiap
$D_i$ termasuk dalam domain $\tilde\mu_i$ dan $D_i=A_i$ untuk semua kecuali
sejumlah berhingga $i$.   Sekarang, untuk setiap $i$ kita dapat menemukan $C_i\in\Sigma_i$ sedemikian
sehingga $D_i=C_i\cap A_i$ dan $\mu_iC_i=\tilde\mu_iD_i$, dan kita dapat mengandaikan
bahwa $C_i=X_i$ kapan pun $D_i=A_i$.   Dalam hal ini $C=\prod_{i\in
I}C_i\in\Lambda$ dan

\Centerline{$\lambda C=\prod_{i\in I}\mu_iC_i=\prod_{i\in
I}\tilde\mu_iD_i=\tilde\lambda D$.}

\noindent Oleh karena itu

\Centerline{$\lambda_A\phi^{-1}[D]=\lambda_A(A\cap C)\le\lambda C
=\tilde\lambda D$.}

\noindent Sekali lagi menurut 254G, $\phi$ bersifat \imp\ dalam perwujudan ini, yaitu
$\lambda_AV$ terdefinisi dan sama dengan $\tilde\lambda V$ untuk setiap
$V\in\tilde\Lambda$.   Dengan menggabungkan ini dengan (i), kita
memperoleh $\lambda_A=\tilde\lambda$, sebagaimana dinyatakan.

\medskip

{\bf (b)} Untuk setiap $i\in I$, pilih himpunan $E_i\in\Sigma_i$ sedemikian sehingga
$A_i\subseteq E_i$ dan $\mu_iE_i=\mu^*_iA_i$;  lakukan ini sedemikian rupa
sehingga $E_i=X_i$ kapan pun $\mu^*_iA_i=1$.   Tetapkan
$B_i=A_i\cup(X_i\setminus E_i)$, sehingga $\mu^*_iB_i=1$ untuk setiap $i$ (jika
$F\in\Sigma_i$ dan
$F\supseteq B_i$ maka $F\cap E_i\supseteq A_i$, sehingga

\Centerline{$\mu_iF=\mu_i(F\cap E_i)+\mu_i(F\setminus E_i)
=\mu_i E_i+\mu_i(X_i\setminus E_i)=1$.)}

\noindent Menurut (a), kita dapat mengidentifikasikan ukuran subruang $\lambda_B$ pada
$B=\prod_{i\in I}B_i$ dengan produk dari ukuran-ukuran subruang
$\tilde\mu_i$ pada $B_i$.   Khususnya, $\lambda^*B=\lambda_BB=1$.
Sekarang $A_i=B_i\cap E_i$ sehingga (dengan menuliskan $A=\prod_{i\in I}A_i$),
$A=B\cap\prod_{i\in I}E_i$.

Jika $\prod_{i\in I}\mu_i^*A_i=0$, maka untuk setiap $\epsilon>0$ terdapat
himpunan berhingga $J\subseteq I$ sedemikian sehingga $\prod_{i\in J}\mu_i^*A_i\le\epsilon$;
akibatnya (dengan menggunakan 254Fb)

\Centerline{$\lambda^*A\le\lambda\{x:x(i)\in E_i$ untuk setiap
$i\in J\}=\prod_{i\in J}\mu_iE_i\le\epsilon$.}

\noindent Karena $\epsilon$ sebarang, $\lambda^*A=0$.  Jika
$\prod_{i\in I}\mu_i^*A_i>0$, maka untuk setiap $n\in\Bbb N$ himpunan
$\{i:\mu_i^*A_i\le 1-2^{-n}\}$ harus berhingga, sehingga

\Centerline{$J=\{i:\mu_i^*A_i<1\}=\{i:E_i\ne X_i\}$}

\noindent terhitung.   Sekali lagi menurut 254Fb,
yang diterapkan pada $\langle E_i\cap B_i\rangle_{i\in I}$ dalam produk
$\prod_{i\in I}B_i$,

$$\eqalign{\lambda^*(\prod_{i\in I}A_i)
&=\lambda_B(\prod_{i\in I}A_i)
=\lambda_B\{x:x\in B,\,x(i)\in E_i\cap B_i\text{ untuk setiap }i\in J\}\cr
&=\prod_{i\in J}\tilde\mu_{i}(E_i\cap B_i)
=\prod_{i\in I}\mu_i^*A_i,\cr}$$

\noindent sebagaimana diperlukan.
}%end of proof of 254L

\leader{254M}{}\cmmnt{ Sekarang saya beralih ke hasil-hasil dasar yang memungkinkan
kita menggunakan ukuran-ukuran produk ini secara efektif.   Mula-mula, saya memperkenalkan
kosakata untuk
menangani subproduk.}   Misalkan $\langle X_i\rangle_{i\in I}$ suatu
keluarga himpunan, dengan produk $X$.

\spheader 254Ma Untuk $J\subseteq I$, tuliskan $X_J$ untuk
$\prod_{i\in J}X_i$.   Kita mempunyai bijeksi kanonik
$x\mapsto(x\restr J,x\restr I\setminus J):
X\to X_J\times X_{I\setminus J}$.   Terkait dengan bijeksi ini,
kita mempunyai peta $x\mapsto\pi_J(x)=x\restr J:X\to X_J$.   Sekarang saya akan mengatakan
bahwa himpunan $W\subseteq X$ {\bf ditentukan oleh koordinat dalam $J$} jika
terdapat $V\subseteq X_J$ sedemikian sehingga $W=\pi_J^{-1}[V]$;  yaitu, $W$
bersesuaian dengan
$V\times X_{I\setminus J}\subseteq X_J\times X_{I\setminus J}$.

\cmmnt{Mudah dilihat bahwa}

$$\eqalign{W&\text{ ditentukan oleh koordinat dalam }J\cr
&\quad\iff x'\in W\text{ kapan pun }x\in W,\,x'\in X\text{ dan }x'\restr
J=x\restr J\cr
&\quad\iff W=\pi_J^{-1}[\pi_J[W]].\cr}$$

\noindent Akibatnya, jika $W$ ditentukan oleh koordinat dalam $J$,
dan $J\subseteq K\subseteq I$, maka $W$ juga ditentukan oleh koordinat dalam
$K$.   Keluarga $\Cal W_J$ dari subhimpunan-subhimpunan $X$ yang ditentukan oleh koordinat
dalam $J$ tertutup terhadap pengambilan komplemen serta gabungan dan
irisan sebarang.
\prooflet{\Prf\ Jika $W\in\Cal W_J$, maka

\Centerline{$X\setminus W=X\setminus\pi_J^{-1}[\pi_J[W]]
=\pi_J^{-1}[X_J\setminus\pi_J[W]]\in\Cal W_J$.}

\noindent Jika $\Cal V\subseteq\Cal W_J$, maka

\Centerline{$\bigcup\Cal V
=\bigcup_{V\in\Cal V}\pi_J^{-1}[\pi_J[V]]
=\pi_J^{-1}[\bigcup_{V\in\Cal V}\pi_J[V]]
\in\Cal W_J$.  \Qed}
}%end of prooflet

\spheader 254Mb \cmmnt{Akibatnya,}

\Centerline{$\Cal W=\bigcup\{\Cal W_J:J\subseteq I$ dengan J terhitung$\}$\cmmnt{,}}

\noindent\cmmnt{keluarga subhimpunan $X$ yang ditentukan oleh koordinat
dalam suatu himpunan
terhitung, }merupakan aljabar-sigma dari subhimpunan-subhimpunan $X$.
\prooflet{\Prf\ (i) $X$ dan $\emptyset$ ditentukan oleh koordinat dalam
$\emptyset$ (ingat bahwa $X_{\emptyset}$ merupakan himpunan tunggal, dan bahwa
$X=\pi_{\emptyset}^{-1}[X_{\emptyset}]$,
$\emptyset=\pi^{-1}_{\emptyset}[\emptyset]$).   (ii) Jika $W\in\Cal W$,
terdapat himpunan terhitung $J\subseteq I$ sedemikian sehingga $W\in\Cal W_J$;  sekarang

\Centerline{$X\setminus W=\pi_J^{-1}[X_J\setminus\pi_J[W]]\in\Cal
W_J\subseteq\Cal W$.}

\noindent (iii) Jika $\sequencen{W_n}$
merupakan barisan dalam $\Cal W$, maka untuk setiap $n\in\Bbb N$ terdapat
himpunan terhitung $J_n\subseteq I$ sedemikian sehingga $W_n\in\Cal W_{J_n}$.   Sekarang
$J=\bigcup_{n\in\Bbb N}J_n$ merupakan subhimpunan terhitung dari $I$,
dan setiap $W_n$ termasuk dalam $\Cal W_J$, sehingga

\Centerline{$\bigcup_{n\in\Bbb N}W_n\in\Cal W_J\subseteq\Cal W$.  \Qed}
}%end of prooflet

\spheader 254Mc Jika $i\in I$ dan $E\subseteq X_i$, maka
$\{x:x\in X$, $x(i)\in E\}$ ditentukan oleh koordinat tunggal $i$,
sehingga tentu saja termasuk dalam $\Cal W$;  dengan demikian $\Cal W$ harus memuat
$\Tensorhat_{i\in I}\Cal PX_i$.
\cmmnt{{\it A fortiori}, jika $\Sigma_i$ merupakan aljabar-sigma dari
subhimpunan-subhimpunan $X_i$ untuk setiap $i$,
$\Cal W\supseteq\Tensorhat_{i\in I}\Sigma_i$;  yaitu,
setiap anggota $\Tensorhat_{i\in I}\Sigma_i$ ditentukan
oleh koordinat dalam suatu himpunan terhitung.}

\vleader{72pt}{254N}{\bf Teorema} Misalkan
$\langle(X_i,\Sigma_i,\mu_i)\rangle_{i\in I}$ suatu
keluarga ruang probabilitas dan $\langle K_j\rangle_{j\in J}$ suatu
partisi dari $I$.   Untuk setiap $j\in J$, misalkan $\lambda_j$ ukuran produk
pada $Z_j=\prod_{i\in K_j}X_i$, dan tuliskan $\lambda$ untuk
ukuran produk pada $X=\prod_{i\in I}X_i$.   Maka bijeksi alami

\Centerline{$x\mapsto\phi(x)
=\langle x\restr K_j\rangle_{j\in J}:X\to\prod_{j\in J}Z_j$}

\noindent mengidentifikasikan $\lambda$ dengan produk dari keluarga
$\langle\lambda_j\rangle_{j\in J}$.

Khususnya, jika $K\subseteq I$ adalah sebarang himpunan, maka $\lambda$ dapat
diidentifikasikan dengan produk c.l.d. dari ukuran-ukuran produk pada
$\prod_{i\in K}X_i$ dan $\prod_{i\in I\setminus K}X_i$.

\proof{ (Bandingkan 251N.) Tuliskan $Z=\prod_{j\in J}Z_j$ dan
$\tilde\lambda$ untuk ukuran produk pada $Z$;  misalkan $\Lambda$,
$\tilde\Lambda$ masing-masing domain dari $\lambda$ dan $\tilde\lambda$.

\medskip

{\bf (a)} Misalkan $C\subseteq Z$ suatu silinder terukur.   Maka
$\lambda^*\phi^{-1}[C]\le\tilde\lambda C$.   \Prf\ Nyatakan $C$ sebagai
$\prod_{j\in J}C_j$ dengan $C_j\subseteq Z_j$ termasuk dalam domain
$\Lambda_j$ dari $\lambda_j$ untuk setiap $j$.
Tetapkan $L=\{j:C_j\ne Z_j\}$, sehingga $L$ berhingga.   Misalkan $\epsilon>0$.
Untuk setiap $j\in L$, misalkan $\sequencen{C_{jn}}$ suatu barisan silinder terukur
dalam $Z_j=\prod_{i\in K_j}X_i$ sedemikian sehingga
$C_j\subseteq\bigcup_{n\in\Bbb N}C_{jn}$ dan
$\sum_{n=0}^{\infty}\lambda_jC_{jn}\le\lambda_jC_j+\epsilon$.   Nyatakan
setiap $C_{jn}$ sebagai $\prod_{i\in K_j}C_{jni}$ dengan $C_{jni}\in\Sigma_i$
untuk $i\in K_j$ (dan $\{i:C_{jni}\ne X_i\}$ berhingga).

Untuk $f\in\BbbN^L$, tetapkan

\Centerline{$D_f
=\{x:x\in X$, $x(i)\in C_{j,f(j),i}$ kapan pun $j\in L$, $i\in K_j\}$.}

\noindent Karena $\bigcup_{j\in L}\{i:C_{j,f(j),i}\ne X_i\}$ berhingga,
$D_f$ merupakan silinder terukur dalam $X$, dan

\Centerline{$\lambda D_f=\prod_{j\in L}\prod_{i\in K_j}\mu_iC_{j,f(j),i}
=\prod_{j\in L}\lambda_jC_{j,f(j)}$.}

\noindent Selain itu,

\Centerline{$\bigcup\{D_f:f\in\BbbN^L\}\supseteq\phi^{-1}[C]$}

\noindent karena jika $\phi(x)\in C$, maka $\phi(x)(j)\in C_j$ untuk setiap
$j\in L$, sehingga pasti terdapat $f\in\BbbN^L$ sedemikian sehingga $\phi(x)(j)\in
C_{j,f(j)}$ untuk setiap $j\in L$.   Tetapi (karena $\BbbN^L$ terhitung)
ini berarti bahwa

$$\eqalign{\lambda^*\phi^{-1}[C]
&\le\sum_{f\in\BbbN^L}\lambda D_f
=\sum_{f\in \BbbN^L}\prod_{j\in L}\lambda_jC_{j,f(j)}\cr
&=\prod_{j\in L}\sum_{n=0}^{\infty}\lambda_jC_{jn}
\le\prod_{j\in L}(\lambda_jC_j+\epsilon).\cr}$$

\noindent Karena $\epsilon$ sebarang,

\Centerline{$\lambda^*\phi^{-1}[C]\le\prod_{j\in L}\lambda_jC_j
=\tilde\lambda C$.  \Qed}

Menurut 254G, kita memperoleh bahwa $\lambda\phi^{-1}[W]$ terdefinisi dan sama dengan
$\tilde\lambda W$ kapan pun $W\in\tilde\Lambda$.
\medskip

{\bf (b)} Selanjutnya, $\tilde\lambda\phi[D]=\lambda D$ untuk setiap silinder terukur
$D\subseteq X$.   \Prf\ Ini mudah.   Nyatakan $D$ sebagai
$\prod_{i\in I}D_i$ dengan $D_i\in\Sigma_i$ untuk setiap $i\in I$ dan
$\{i:D_i\ne X_i\}$ berhingga.   Maka $\phi[D]=\prod_{j\in J}\tilde
D_j$, dengan $\tilde D_j=\prod_{i\in K_j}D_i$ suatu silinder terukur
untuk setiap $j\in J$.   Karena $\{j:\tilde D_j\ne Z_j\}$ juga harus
berhingga (bahkan, himpunan ini tidak dapat mempunyai lebih banyak anggota daripada himpunan berhingga
$\{i:D_i\ne X_i\}$), $\prod_{j\in J}\tilde D_j$ sendiri merupakan silinder terukur
dalam $Z$, dan

\Centerline{$\tilde\lambda\phi[D]
=\prod_{j\in J}\lambda_j\tilde D_j
=\prod_{j\in J}\prod_{i\in K_j}\mu_iD_i
=\lambda D$.   \Qed}

Dengan menerapkan 254G pada $\phi^{-1}:Z\to X$, diperoleh bahwa
$\tilde\lambda\phi[W]$ terdefinisi dan sama dengan $\lambda W$ untuk setiap
$W\in\Lambda$.   Namun, bersama dengan (a), ini berarti bahwa
untuk sebarang $W\subseteq X$,

\inset{jika $W\in\Lambda$, maka $\phi[W]\in\tilde\Lambda$ dan
$\tilde\lambda\phi[W]=\lambda W$,}

\inset{jika $\phi[W]\in\tilde\Lambda$, maka $W\in\Lambda$ dan $\lambda
W=\tilde\lambda\phi[W]$.}

\noindent Dan tentu saja, inilah tepatnya arti pernyataan bahwa $\phi$
merupakan isomorfisme antara $(X,\Lambda,\lambda)$ dan
$(Z,\tilde\Lambda,\tilde\lambda)$.
}%end of proof of 254N

\leader{254O}{Proposisi} Misalkan
$\langle(X_i,\Sigma_i,\mu_i)\rangle_{i\in I}$ suatu keluarga ruang probabilitas.
Untuk setiap $J\subseteq I$, misalkan
$\lambda_J$ ukuran probabilitas produk pada
$X_J=\prod_{i\in J}X_i$, dan $\Lambda_J$ domainnya;  tuliskan $X=X_I$,
$\lambda=\lambda_I$ dan $\Lambda=\Lambda_I$.   Untuk $x\in X$ dan
$J\subseteq I$, tetapkan $\pi_J(x)=x\restr J\in X_J$.

(a) Untuk setiap $J\subseteq I$, $\lambda_J$ adalah ukuran citra
$\lambda\pi_J^{-1}$\cmmnt{ (234D)};  khususnya,
$\pi_J:X\to X_J$ bersifat \imp\ untuk $\lambda$
dan $\lambda_J$.

(b) Jika $J\subseteq I$ dan $W\in\Lambda$
ditentukan oleh koordinat dalam $J$\cmmnt{ (254M)}, maka
$\lambda_J\pi_J[W]$ terdefinisi dan sama dengan $\lambda W$.   Akibatnya,
terdapat $W_1$, $W_2$ yang termasuk dalam aljabar-sigma dari subhimpunan-subhimpunan
$X$ yang dibangkitkan oleh

\Centerline{$\{\{x:x(i)\in E\}:i\in J$, $E\in\Sigma_i\}$}

\noindent sedemikian sehingga $W_1\subseteq W\subseteq W_2$ dan
$\lambda(W_2\setminus W_1)=0$.

(c) Untuk setiap $W\in\Lambda$, kita dapat menemukan suatu
himpunan terhitung $J$ dan $W_1$, $W_2\in\Lambda$, keduanya
ditentukan oleh koordinat dalam $J$, sedemikian sehingga
$W_1\subseteq W\subseteq W_2$ dan $\lambda(W_2\setminus W_1)=0$.

(d) Untuk setiap $W\in\Lambda$, terdapat
himpunan terhitung $J\subseteq I$ sedemikian sehingga $\pi_J[W]\in\Lambda_J$ dan
$\lambda_J\pi_J[W]=\lambda W$;  sehingga
$W'=\pi_J^{-1}[\pi_J[W]]$ termasuk dalam $\Lambda$, dan
$\lambda(W'\setminus W)=0$.

\proof{{\bf (a)(i)} Menurut 254N, kita dapat mengidentifikasikan $\lambda$ dengan
produk dari $\lambda_J$ dan $\lambda_{I\setminus J}$ pada
$X_J\times X_{I\setminus J}$.
Sekarang $\pi_J^{-1}[E]\subseteq X$ bersesuaian dengan
$E\times X_{I\setminus J}\subseteq X_J\times X_{I\setminus J}$, sehingga

\Centerline{$\lambda(\pi_J^{-1}[E])
=\lambda_JE\cdot\lambda_{I\setminus J}X_{I\setminus J}=\lambda_JE$,}

\noindent menurut 251E atau 251Ia, kapan pun $E\in\Lambda_J$.   Ini menunjukkan bahwa
$\pi_J$ bersifat \imp.

\medskip

\quad{\bf (ii)} Untuk melihat bahwa $\lambda_J$ benar-benar merupakan ukuran citra,
andaikan bahwa $E\subseteq X_J$ sedemikian sehingga $\pi_J^{-1}[E]\in\Lambda$.
Dengan mengidentifikasikan $\pi^{-1}_J[E]$ dengan $E\times X_{I\setminus J}$, seperti sebelumnya,
kita mengandaikan bahwa $E\times X_{I\setminus J}$ diukur oleh
ukuran produk pada $X_J\times X_{I\setminus J}$.   Namun, ini berarti bahwa
untuk $\lambda_{I\setminus J}$-hampir setiap
$z\in X_{I\setminus J}$, $E_z=\{y:(y,z)\in E\times X_{I\setminus J}\}$
termasuk dalam $\Lambda_J$ (252D(ii), karena $\lambda_J$ lengkap).
Karena $E_z=E$ untuk setiap $z$, $E$ sendiri termasuk dalam $\Lambda_J$, sebagaimana
dinyatakan.

\medskip

{\bf (b)} Jika $W\in\Lambda$ ditentukan oleh koordinat
dalam $J$, tetapkan $H=\pi_J[W]$;  maka $\pi_J^{-1}[H]=W$, sehingga $H\in\Lambda_J$ menurut
(a) tepat di atas.   Menurut 254Ff, terdapat $H_1$,
$H_2\in\Tensorhat_{i\in J}\Sigma_i$ sedemikian sehingga
$H_1\subseteq H\subseteq H_2$ dan $\lambda_J(H_2\setminus H_1)=0$.

Misalkan $\Tau_J$ aljabar-sigma dari subhimpunan-subhimpunan $X$ yang dibangkitkan oleh himpunan-himpunan
berbentuk $\{x:x(i)\in E\}$ dengan $i\in J$ dan $E\in\Sigma_i$.
Tinjau $\Tau'_J=\{G:G\subseteq X_J$, $\pi_J^{-1}[G]\in\Tau_J\}$.
Keluarga ini
merupakan aljabar-sigma dari subhimpunan-subhimpunan $X_J$, dan memuat $\{y:y\in
X_J$, $y(i)\in E\}$ kapan pun $i\in J$, $E\in\Sigma_i$ (karena

\Centerline{$\pi_J^{-1}[\{y:y\in X_J$, $y(i)\in E\}]=\{x:x\in X$,
$x(i)\in
E\}$}

\noindent kapan pun $i\in J$, $E\subseteq X_i$).   Jadi $\Tau'_J$ harus
memuat $\Tensorhat_{i\in J}\Sigma_i$.   Khususnya, $H_1$ dan $H_2$
keduanya termasuk dalam $\Tau'_J$, yaitu $W_k=\pi_J^{-1}[H_k]$ termasuk dalam
$\Tau_J$ untuk kedua nilai $k$.   Tentu saja $W_1\subseteq W\subseteq W_2$,
karena $H_1\subseteq H\subseteq H_2$, dan

\Centerline{$\lambda(W_2\setminus W_1)=\lambda_J(H_2\setminus H_1)=0$,}
\noindent sebagaimana diperlukan.

\medskip

{\bf (c)} Sekarang ambil sebarang $W\in\Lambda$.   Menurut
254Ff, terdapat $W_1$ dan $W_2\in\Tensorhat_{i\in I}\Sigma_i$ sedemikian sehingga
$W_1\subseteq W\subseteq W_2$ dan $\lambda(W_2\setminus W_1)=0$.   Menurut
254Mc, terdapat himpunan-himpunan terhitung $J_1$, $J_2\subseteq I$ sedemikian sehingga, untuk
setiap $k$, $W_k$ ditentukan oleh koordinat dalam $J_k$.   Dengan menetapkan
$J=J_1\cup J_2$, $J$ merupakan subhimpunan terhitung dari $I$ dan $W_1$ maupun
$W_2$ ditentukan oleh koordinat dalam $J$.

\medskip

{\bf (d)} Dengan melanjutkan argumen dari (c), $\pi_J[W_1]$,
$\pi_J[W_2]\in\Lambda_J$, menurut (b), dan
$\lambda_J(\pi_J[W_2]\setminus\pi_J[W_1])=0$.   Karena
$\pi_J[W_1]\subseteq\pi_J[W]\subseteq\pi_J[W_2]$, diperoleh bahwa
$\pi_J[W]\in\Lambda_J$, dengan
$\lambda_J\pi_J[W]=\lambda_J\pi_J[W_2]$;  sehingga, dengan menetapkan
$W'=\pi_J^{-1}[\pi_J[W]]$, $W'\in\Lambda$, dan

\Centerline{$\lambda W'=\lambda_J\pi_J[W]=\lambda_J\pi_J[W_2]
=\lambda\pi_J^{-1}[\pi_J[W_2]]=\lambda W_2=\lambda W$.}
}%end of proof of 254O

\leader{254P}{Proposisi} Misalkan
$\langle(X_i,\Sigma_i,\mu_i)\rangle_{i\in I}$
suatu keluarga ruang probabilitas, dan untuk setiap $J\subseteq I$ misalkan
$\lambda_J$ ukuran probabilitas produk pada
$X_J=\prod_{i\in J}X_i$, dan $\Lambda_J$ domainnya;  tuliskan $X=X_I$,
$\Lambda=\Lambda_I$ dan $\lambda=\lambda_I$.   Untuk $x\in X$ dan
$J\subseteq I$, tetapkan $\pi_J(x)=x\restr J\in X_J$.

(a) Jika $J\subseteq I$ dan $g$ suatu fungsi bernilai real yang didefinisikan pada suatu
subhimpunan $X_J$, maka $g$ terukur terhadap $\Lambda_J$ jika dan hanya jika $g\pi_J$
terukur terhadap $\Lambda$.

(b) Kapan pun $f$ suatu fungsi bernilai real yang terukur terhadap $\Lambda$ dan didefinisikan
pada suatu subhimpunan yang koterabaikan terhadap $\lambda$ di dalam $X$, kita dapat menemukan suatu himpunan terhitung
$J\subseteq I$ dan suatu fungsi yang terukur terhadap $\Lambda_J$, yakni $g$, yang didefinisikan pada suatu
subhimpunan yang koterabaikan terhadap $\lambda_J$ di dalam $X_J$ sedemikian sehingga $f$ memperluas $g\pi_J$.

\proof{{\bf (a)(i)} Jika $g$ terukur terhadap $\Lambda_J$ dan $a\in\Bbb R$,
terdapat $H\in\Lambda_J$ sedemikian sehingga $\{y:y\in\dom g$, $g(y)\ge
a\}=H\cap\dom g$.   Sekarang $\pi_J^{-1}[H]\in\Lambda$, menurut 254Oa, dan
$\{x:x\in\dom g\pi_J$, $g\pi_J(x)\ge a\}=\pi_J^{-1}[H]\cap\dom g\pi_J$.
Jadi
$g\pi_J$ terukur terhadap $\Lambda$.

\medskip

\quad{\bf (ii)} Jika $g\pi_J$ terukur terhadap $\Lambda$ dan $a\in\Bbb R$,
maka terdapat $W\in\Lambda$ sedemikian sehingga
$\{x:g\pi_J(x)\ge a\}=W\cap\dom g\pi_J$.   Seperti dalam bukti 254Oa, kita
dapat mengidentifikasikan $\lambda$ dengan
produk dari $\lambda_J$ dan $\lambda_{I\setminus J}$, dan 252D(ii) memberi tahu
kita bahwa, jika kita mengidentifikasikan $W$ dengan subhimpunan yang bersesuaian dari
$X_J\times X_{I\setminus J}$, terdapat setidaknya satu
$z\in X_{I\setminus J}$ sedemikian
sehingga $W_z=\{y:y\in X_J$, $(y,z)\in W\}$ termasuk dalam $\Lambda_J$.   Namun,
karena (dengan konvensi ini) $g\pi_J(y,z)=g(y)$ untuk setiap $y\in X_J$, kita
melihat bahwa $\{y:y\in\dom g$, $g(y)\ge a\}=W_z\cap\dom g$.   Karena $a$
sebarang, $g$ terukur terhadap $\Lambda_J$.

\medskip

{\bf (b)} Untuk bilangan-bilangan rasional $q$, tetapkan
$W_q=\{x:x\in\dom f$, $f(x)\ge q\}$.   Menurut 254Oc, untuk setiap $q$ kita dapat menemukan
himpunan terhitung
$J_q\subseteq I$ dan himpunan-himpunan $W'_q$, $W_q''$, yang keduanya ditentukan oleh
koordinat
dalam $J_q$, sedemikian sehingga $W'_q\subseteq W_q\subseteq W_q''$ dan
$\lambda(W_q''\setminus W_q')=0$.   Tetapkan $J=\bigcup_{q\in\Bbb Q}J_q$,
$V=X\setminus\bigcup_{q\in\Bbb Q}(W_q''\setminus W_q')$;  maka $J$ merupakan
subhimpunan terhitung dari $I$ dan $V$ merupakan subhimpunan koterabaikan dari $X$;
selain itu, $V$ ditentukan oleh koordinat dalam $J$ karena semua
$W_q'$, $W_q''$ demikian.

Untuk setiap $q\in\Bbb Q$, $W_q\cap V=W_q'\cap V$, karena
$V\cap(W_q\setminus W_q')\subseteq V\cap(W_q''\setminus
W_q')=\emptyset$;  jadi $W_q\cap V$ ditentukan oleh koordinat dalam $J$.
Akibatnya, $V\cap\dom f=\bigcup_{q\in\Bbb Q}V\cap W_q$ juga
ditentukan oleh koordinat dalam $J$.   Selain itu,

\Centerline{$\{x:x\in V\cap\dom f$, $f(x)\ge a\}
=\bigcap_{q\le a}V\cap W_q$}

\noindent ditentukan oleh koordinat dalam $J$.   Artinya,
jika $x$, $x'\in V$ dan $\pi_Jx=\pi_Jx'$, maka $x\in\dom f$ jika dan hanya jika $x'\in\dom
f$, dan dalam hal ini $f(x)=f(x')$.   Dengan menetapkan $H=\pi_J[V\cap\dom f]$, kita
mempunyai $\pi_J^{-1}[H]=V\cap\dom f$, suatu subhimpunan koterabaikan dari $X$, sehingga
(karena $\lambda_J=\lambda\pi_J^{-1}$) $H$ koterabaikan dalam $X_J$.
Selain itu, untuk
$y\in H$, $f(x)=f(x')$ kapan pun $\pi_Jx=\pi_Jx'=y$, sehingga terdapat suatu
fungsi $g:H\to\Bbb R$ yang didefinisikan dengan menetapkan bahwa $g\pi_J(x)=f(x)$ kapan pun
$x\in V\cap\dom f$.   Dengan demikian, $g$ didefinisikan hampir di mana-mana dalam $X_J$ dan
$f$ memperluas $g\pi_J$.   Akhirnya, untuk sebarang $a\in\Bbb R$,

\Centerline{$\pi_J^{-1}[\{y:g(y)\ge a\}]
=\{x:x\in V\cap\dom f$, $f(x)\ge a\}\in\Lambda$;}

\noindent menurut 254Oa, $\{y:g(y)\ge a\}\in\Lambda_J$;  karena $a$ sebarang,
$g$ terukur.
}%end of proof of 254P

\leader{254Q}{Proposisi} Misalkan $\langle(X_i,\Sigma_i,\mu_i)\rangle_{i\in
I}$ suatu keluarga ruang probabilitas, dan untuk setiap $J\subseteq I$ misalkan
$\lambda_J$ ukuran probabilitas produk pada $X_J=\prod_{i\in
J}X_i$;  tuliskan $X=X_I$, $\lambda=\lambda_I$.   Untuk $x\in X$, $J\subseteq
I$, tetapkan $\pi_J(x)=x\restr J\in X_J$.

(a) Misalkan $\eusm S$ subruang linear dari $\Bbb R^X$ yang direntang oleh
$\{\chi C:C\subseteq X$ adalah silinder terukur$\}$.   Maka untuk setiap
fungsi bernilai real yang dapat diintegralkan terhadap $\lambda$, yakni $f$, dan setiap $\epsilon>0$,
terdapat $g\in\eusm S$ sedemikian sehingga $\int|f-g|d\lambda\le\epsilon$.

(b) Kapan pun $J\subseteq I$ dan $g$ suatu fungsi bernilai real yang didefinisikan pada
suatu subhimpunan $X_J$, maka $\int g\,d\lambda_J=\int g\pi_Jd\lambda$ jika
salah satu dari kedua integral tersebut terdefinisi dalam $[-\infty,\infty]$.

(c) Kapan pun $f$ suatu fungsi bernilai real yang dapat diintegralkan terhadap $\lambda$, kita dapat
menemukan himpunan terhitung $J\subseteq I$ dan suatu
fungsi yang dapat diintegralkan terhadap $\lambda_J$, yakni $g$, sedemikian sehingga $f$ memperluas $g\pi_J$.

\proof{{\bf (a)(i)} Tuliskan $\overline{\eusm S}$ untuk himpunan fungsi-fungsi
$f$ yang memenuhi pernyataan tersebut, yaitu sedemikian sehingga untuk setiap $\epsilon>0$
terdapat $g\in\eusm S$ sedemikian sehingga $\int|f-g|\le\epsilon$.   Maka
$f_1+f_2$ dan $cf_1\in\overline{\eusm S}$ kapan pun $f_1$,
$f_2\in\overline{\eusm S}$.   \Prf\ Diberikan $\epsilon>0$, terdapat $g_1$,
$g_2\in\eusm S$ sedemikian sehingga $\int|f_1-g_1|\le\bover{\epsilon}{2+|c|}$,
$\int|f_2-g_2|\le\bover{\epsilon}2$;  sekarang $g_1+g_2$, $cg_1\in\eusm S$
dan $\int|(f_1+f_2)-(g_1+g_2)|\le\epsilon$,
$\int|cf_1-cg_1|\le\epsilon$.\ \QeD\  Selain itu, tentu saja,
$f\in\overline{\eusm S}$ kapan pun $f_0\in\overline{\eusm S}$ dan
$f\eae f_0$.


\medskip

\quad{\bf (ii)} Tuliskan $\Cal W$ untuk
$\{W:W\subseteq X$, $\chi W\in\overline{\eusm S}\}$, dan $\Cal C$ untuk
keluarga
silinder terukur
dalam $X$.   Maka jelas dari definisi dalam 254A bahwa
$C\cap C'\in\Cal C$ untuk semua $C$, $C'\in\Cal C$, dan tentu saja
$C\in\Cal W$ untuk setiap $C\in\Cal C$, karena $\chi C\in\eusm S$.   Selanjutnya,
$W\setminus V\in\Cal W$ kapan pun $W$, $V\in\Cal W$ dan $V\subseteq W$,
karena dalam hal ini $\chi(W\setminus V)=\chi W-\chi V$.   Ketiga,
$\bigcup_{n\in\Bbb N}W_n\in\Cal W$ untuk setiap barisan tak-menurun
$\sequencen{W_n}$ dalam $\Cal W$.   \Prf\ Tetapkan $W=\bigcup_{n\in\Bbb N}W_n$.
Diberikan $\epsilon>0$, terdapat $n\in\Bbb N$ sedemikian sehingga
$\lambda(W\setminus W_n)\le\bover{\epsilon}2$.   Sekarang terdapat
$g\in\eusm S$ sedemikian sehingga $\int|\chi W_n-g|\le\bover{\epsilon}2$, sehingga
$\int|\chi W-g|\le\epsilon$.\ \QeD\  Jadi $\Cal W$ merupakan kelas Dynkin dari
subhimpunan-subhimpunan $X$.

Menurut Teorema Kelas Monoton (136B), $\Cal W$ harus memuat
aljabar-sigma dari subhimpunan-subhimpunan $X$ yang dibangkitkan oleh $\Cal C$, yaitu
$\Tensorhat_{i\in I}\Sigma_i$.   Namun, ini berarti bahwa $\Cal W$ memuat
setiap subhimpunan terukur dari $X$, karena menurut 254Ff setiap himpunan terukur
berbeda dengan suatu anggota $\Tensorhat_{i\in
I}\Sigma_i$ hanya pada suatu himpunan terabaikan.

\medskip

\quad{\bf (iii)} Dengan demikian, $\overline{\eusm S}$ memuat fungsi indikator
dari sebarang subhimpunan terukur dari $X$.   Karena tertutup terhadap
penjumlahan dan perkalian skalar, himpunan ini memuat semua fungsi sederhana.
Namun, ini berarti bahwa himpunan tersebut harus memuat semua fungsi terintegralkan.   \Prf\ Jika
$f$ suatu fungsi bernilai real yang terintegralkan pada $X$, dan
$\epsilon>0$, terdapat fungsi sederhana $h:X\to\Bbb R$ sedemikian sehingga
$\int|f-h|\le\bover{\epsilon}2$ (242M), dan sekarang terdapat $g\in\eusm S$
sedemikian sehingga $\int|h-g|\le\bover{\epsilon}2$, sehingga
$\int|f-g|\le\epsilon$.\ \Qed

Ini membuktikan bagian (a) dari proposisi.

\medskip

{\bf (b)} Gabungkan 254Oa dan 235J.

\medskip

{\bf (c)} Menurut 254Pb, terdapat himpunan terhitung $J\subseteq I$ dan suatu
fungsi bernilai real $g$ yang didefinisikan pada suatu subhimpunan koterabaikan dari $X_J$ sedemikian
sehingga $f$ memperluas $g\pi_J$.   Sekarang $\dom(g\pi_J)=\pi_J^{-1}[\dom g]$
merupakan subhimpunan koterabaikan, sehingga $f\eae g\pi_J$ dan $g\pi_J$ dapat diintegralkan terhadap $\lambda$.
Menurut (b), $g$ dapat diintegralkan terhadap $\lambda_J$.
}%end of proof of 254Q

\leader{254R}{Ekspektasi
bersyarat \dvrocolon{lagi}}\cmmnt{ Dengan menggabungkan gagasan-gagasan 253H
dengan uraian di atas, kita memperoleh beberapa hasil yang penting
bukan hanya karena penerapan langsungnya, melainkan juga karena memperjelas
struktur-struktur di sini.

\medskip

\noindent}{\bf Teorema} Misalkan
$\langle(X_i,\Sigma_i,\mu_i)\rangle_{i\in I}$ suatu keluarga ruang
probabilitas dengan produk
$(X,\Lambda,\lambda)$. Untuk $J\subseteq I$, misalkan
$\Lambda_J\subseteq\Lambda$ subaljabar-sigma dari himpunan-himpunan yang
ditentukan oleh koordinat dalam $J$\cmmnt{ (254Ma)}. Maka kita dapat memandang
$L^0(\lambda\restr\Lambda_J)$ sebagai subruang dari
$L^0(\lambda)$\cmmnt{ (242Jh)}. Misalkan
$P_J:L^1(\lambda)\to L^1(\lambda\restr\Lambda_J)\subseteq L^1(\lambda)$
merupakan operator
ekspektasi bersyarat yang bersesuaian\cmmnt{ (242Jd)}. Maka

(a) untuk sembarang $J$, $K\subseteq I$, $P_{K\cap J}=P_KP_J$;

(b) untuk sembarang $u\in L^1(\lambda)$, terdapat himpunan terhitung
$J^*\subseteq I$ sedemikian sehingga $P_Ju=u$ jika dan hanya jika $J\supseteq J^*$;

(c) untuk sembarang $u\in L^0(\lambda)$, terdapat himpunan terkecil yang tunggal
$J^*\subseteq I$ sedemikian sehingga $u\in L^0(\lambda\restr\Lambda_{J^*})$, dan
$J^*$ ini terhitung;

(d) untuk sembarang $W\in\Lambda$, terdapat himpunan terkecil yang tunggal
$J^*\subseteq I$
sedemikian sehingga $W\symmdiff W'$ terabaikan untuk suatu $W'\in\Lambda_{J^*}$,
dan $J^*$ ini terhitung;

(e) untuk sembarang fungsi bernilai real yang $\Lambda$-terukur $f:X\to\Bbb R$,
terdapat himpunan terkecil yang tunggal $J^*\subseteq I$ sedemikian sehingga $f$ sama
hampir di mana-mana dengan suatu fungsi $\Lambda_{J^*}$-terukur, dan $J^*$ ini
terhitung.

\proof{ Untuk $J\subseteq I$, tuliskan $X_J=\prod_{i\in J}X_i$, misalkan
$\lambda_J$ ukuran produk pada $X_J$, dan tetapkan
$\pi_J(x)=x\restr J$ untuk $x\in X$. Tuliskan $L^0_J$ untuk
$L^0(\lambda\restr\Lambda_J)$, yang dipandang sebagai subhimpunan dari $L^0=L^0_I$, dan
$L^1_J$ untuk $L^1(\lambda\restr\Lambda_J)=L^1(\lambda)\cap L^0_J$, seperti dalam
242Jb; dengan demikian $L^1_J$ adalah himpunan nilai proyeksi $P_J$.

\medskip

{\bf (a)(i)} Misalkan $C\subseteq X$ suatu silinder terukur, yang dinyatakan sebagai
$\prod_{i\in I}C_i$ dengan $C_i\in\Sigma_i$ untuk setiap $i$ dan
$L=\{i:C_i\ne X_i\}$ berhingga. Tetapkan

\Centerline{$C'_i=C_i$ untuk $i\in J$, $X_i$ untuk $i\in I\setminus J$,
\quad$C'=\prod_{i\in I}C'_i$,
\quad$\alpha=\prod_{i\in I\setminus J}\mu_iC_i$.}

\noindent Maka $\alpha\chi C'$ merupakan suatu ekspektasi
bersyarat dari $\chi C$ pada $\Lambda_J$. \Prf\ Menurut
254N, kita dapat mengidentifikasi $\lambda$ dengan produk dari $\lambda_J$ dan
$\lambda_{I\setminus J}$. Ini mengidentifikasi $\Lambda_J$ dengan
$\{E\times X_{I\setminus J}:E\in\dom\lambda_J\}$. Menurut 253H, kita mempunyai
ekspektasi bersyarat $g$ dari $\chi C$ yang didefinisikan dengan menetapkan

\Centerline{$g(y,z)=\int\chi C(y,t)\lambda_{I\setminus J}(dt)$}

\noindent untuk $y\in X_J$, $z\in X_{I\setminus J}$. Tetapi $C$
diidentifikasi dengan $C_J\times C_{I\setminus J}$, dengan
$C_J=\prod_{i\in J}C_i$, sehingga $g(y,z)=0$ jika $y\notin C_J$ dan
jika tidak, nilainya adalah
$\lambda_{I\setminus J}C_{I\setminus J}=\alpha$. Jadi
$g=\alpha\chi(C_J\times X_{I\setminus J})$. Tetapi identifikasi
antara $X_J\times X_{I\setminus J}$ dan $X$ memadankan
$C_J\times X_{I\setminus J}$ dengan $C'$, sebagaimana diuraikan di atas. Jadi $g$
menjadi
teridentifikasi dengan $\alpha\chi C'$, dan $\alpha\chi C'$ merupakan ekspektasi
bersyarat dari $\chi C$.\ \Qed

\medskip

\quad{\bf (ii)} Selanjutnya, dengan menetapkan

\Centerline{$C''_i=C'_i$ untuk $i\in K$, $X_i$ untuk $i\in I\setminus K$,
\quad$C''=\prod_{i\in I}C''_i$,}

\Centerline{$\beta=\prod_{i\in I\setminus K}\mu_iC'_i
=\prod_{i\in I\setminus(J\cup K)}\mu_iC_i$,}

\noindent argumen yang sama menunjukkan bahwa $\beta\chi C''$ merupakan ekspektasi
bersyarat dari $\chi C'$ pada $\Lambda_K$. Jadi kita mempunyai

\Centerline{$P_KP_J(\chi C)^{\ssbullet}
=\beta\alpha(\chi C'')^{\ssbullet}$.}

\noindent Tetapi jika kita memperhatikan $\beta\alpha$, ini hanyalah
$\prod_{i\in I\setminus(K\cap J)}\mu_iC_i$, sedangkan $C''_i=C_i$ jika
$i\in K\cap J$,
$X_i$ untuk $i$ lainnya. Jadi $\beta\alpha\chi C''$ merupakan ekspektasi
bersyarat dari $\chi C$ pada $\Lambda_{K\cap J}$, dan

\Centerline{$P_KP_J(\chi C)^{\ssbullet}
=P_{K\cap J}(\chi C)^{\ssbullet}$.}

\medskip

\quad{\bf (iii)} Dengan demikian kita melihat bahwa operator-operator $P_KP_J$, $P_{K\cap J}$
berimpit pada unsur-unsur berbentuk $\chi C^{\ssbullet}$, dengan $C$ suatu
silinder terukur. Karena keduanya linear, keduanya berimpit pada
kombinasi linear unsur-unsur tersebut, yaitu $P_KP_Jv=P_{K\cap J}v$ kapan pun
$v=g^{\ssbullet}$ untuk suatu $g$ dalam ruang $\eusm S$ dari 254Q. Tetapi jika
$u\in L^1(\lambda)$ dan $\epsilon>0$, terdapat fungsi yang dapat diintegralkan terhadap
$\lambda$, yaitu $f$, sedemikian sehingga $f^{\ssbullet}=u$, dan terdapat $g\in\eusm S$
sedemikian sehingga $\int|f-g|\le\epsilon$ (254Qa), sehingga
$\|u-v\|_1\le\epsilon$, dengan $v=g^{\ssbullet}$. Karena $P_J$, $P_K$
dan $P_{K\cap J}$ semuanya merupakan operator linear bernorma $1$,

\Centerline{$\|P_KP_Ju-P_{K\cap J}u\|_1
\le2\|u-v\|_1+\|P_KP_Jv-P_{K\cap J}v\|_1\le 2\epsilon$.}

\noindent Karena $\epsilon$ sembarang, $P_KP_Ju=P_{K\cap J}u$; karena $u$
sembarang, $P_KP_J=P_{K\cap J}$.

\medskip

{\bf (b)} Ambil $u\in L^1(\lambda)$. Misalkan $\Cal J$ keluarga semua
subhimpunan $J$ dari $I$ sedemikian sehingga $P_Ju=u$. Menurut (a), $J\cap K\in\Cal J$ untuk
semua $J$, $K\in\Cal J$. Selanjutnya, $\Cal J$ memuat suatu himpunan terhitung $J_0$.
\Prf\ Misalkan $f$ suatu fungsi yang dapat diintegralkan terhadap $\lambda$ sedemikian sehingga
$f^{\ssbullet}=u$. Menurut 254Qc, kita dapat menemukan himpunan terhitung
$J_0\subseteq I$ dan fungsi yang dapat diintegralkan terhadap $\lambda_{J_0}$, yaitu $g$, sedemikian sehingga
$f\eae g\pi_{J_0}$. Sekarang $g\pi_{J_0}$ bersifat $\Lambda_{J_0}$-terukur
dan $u=(g\pi_{J_0})^{\ssbullet}$ termasuk dalam $L^1_{J_0}$, sehingga
$J_0\in\Cal J$.\ \Qed

Tuliskan $J^*=\bigcap\Cal J$, sehingga $J^*\subseteq J_0$ terhitung.
Maka $J^*\in\Cal J$. \Prf\ Misalkan $\epsilon>0$. Seperti dalam bukti (a)
di atas, terdapat $g\in\eusm S$ sedemikian sehingga $\|u-v\|_1\le\epsilon$, dengan
$v=g^{\ssbullet}$. Tetapi karena $g$ merupakan kombinasi linear berhingga dari
fungsi-fungsi indikator silinder terukur, yang masing-masing ditentukan oleh
koordinat-koordinat dalam suatu
himpunan berhingga, terdapat $K\subseteq I$ berhingga sedemikian sehingga
$g$ bersifat $\Lambda_K$-terukur, sehingga $P_Kv=v$.
Karena $K$ berhingga, harus ada $J_1,\ldots,J_n\in\Cal J$ sedemikian sehingga
$J^*\cap K=\bigcap_{1\le i\le n}J_i\cap K$; tetapi karena $\Cal J$ tertutup
terhadap irisan berhingga, $J=J_1\cap\ldots\cap J_n\in\Cal J$, dan
$J^*\cap K=J\cap K$.

Sekarang kita mempunyai

\Centerline{$P_{J^*}v=P_{J^*}P_Kv=P_{J^*\cap K}v=P_{J\cap K}v=P_JP_Kv
=P_Jv$,}

\noindent dengan menggunakan (a) dua kali. Karena $P_J$ maupun $P_{J^*}$ bernorma
$1$,

$$\eqalign{\|P_{J^*}u-u\|_1
&\le\|P_{J^*}u-P_{J^*}v\|_1
  +\|P_{J^*}v-P_Jv\|_1
  +\|P_Jv-P_Ju\|_1
  +\|P_Ju-u\|_1\cr
&\le\|u-v\|_1+0+\|u-v\|_1+0
\le 2\epsilon.\cr}$$

\noindent Karena $\epsilon$ sembarang, $P_{J^*}u=u$ dan
$J^*\in\Cal J$.\ \Qed

Sekarang, untuk sembarang $J\subseteq I$,

$$\eqalign{P_Ju=u
&\Longrightarrow J\in\Cal J
\Longrightarrow J\supseteq J^*\cr
&\Longrightarrow P_Ju=P_JP_{J^*}u=P_{J\cap J^*}u=P_{J^*}u=u.\cr}$$

\noindent Jadi $J^*$ mempunyai sifat-sifat yang diminta.

\medskip

{\bf (c)} Tetapkan $e=(\chi X)^{\ssbullet}$, $u_n=(-ne)\vee(u\wedge ne)$ untuk
setiap $n\in\Bbb N$. Maka, untuk sembarang $J\subseteq I$, $u\in L^0_J$ jika dan
hanya jika $u_n\in L^0_J$ untuk setiap $n$. \Prf\ ($\alpha$) Jika $u\in L^0_J$, maka
$u$ dapat dinyatakan sebagai $f^{\ssbullet}$ untuk suatu fungsi yang $\Lambda_J$-terukur
$f$; sekarang $f_n=(-n\chi X)\vee(f\wedge n\chi X)$ bersifat
$\Lambda_J$-terukur, sehingga $u_n=f_n^{\ssbullet}\in L^0_J$ untuk setiap
$n$. ($\beta$) Jika $u_n\in L^0_J$ untuk setiap $n$, maka untuk setiap $n$ kita
dapat menemukan fungsi $\Lambda_J$-terukur $f_n$ sedemikian sehingga
$f_n^{\ssbullet}=u_n$. Tetapi juga terdapat fungsi $\Lambda$-terukur
$f$ sedemikian sehingga $u=f^{\ssbullet}$, dan kita harus mempunyai
$f_n\eae(-n\chi X)\vee(f\wedge n\chi X)$ untuk setiap $n$, sehingga
$f\eae\lim_{n\to\infty}f_n$ dan
$u=(\lim_{n\to\infty}f_n)^{\ssbullet}$. Karena $\lim_{n\to\infty}f_n$
bersifat $\Lambda_J$-terukur dan didefinisikan pada suatu
himpunan koterabaikan terhadap $\lambda\restr\Lambda_J$, maka $u\in L^0_J$.\ \Qed

Karena setiap $u_n$ termasuk dalam $L^1$, kita tahu bahwa

\Centerline{$u_n\in L^0_J\iff u_n\in L^1_J\iff P_Ju_n=u_n$.}

\noindent Menurut (b), untuk setiap $n$ terdapat $J_n^*$ terhitung sedemikian sehingga
$P_{J}u_n=u_n$ jika dan hanya jika $J\supseteq J_n^*$. Jadi kita melihat bahwa $u\in L^0_J$ jika dan
hanya jika $J\supseteq J_n^*$ untuk setiap $n$, yaitu,
$J\supseteq\bigcup_{n\in\Bbb N}J_n^*$. Dengan demikian
$J^*=\bigcup_{n\in\Bbb N}J_n^*$ mempunyai sifat yang dinyatakan.

\medskip

{\bf (d)} Dengan menerapkan (c) pada $u=(\chi W)^{\ssbullet}$, kita memperoleh
$J^*$ (terhitung) terkecil yang tunggal sedemikian sehingga $u\in L^0_{J^*}$. Tetapi jika
$J\subseteq I$, maka terdapat $W'\in\Lambda_J$ sedemikian sehingga
$W'\symmdiff W$ terabaikan jika dan hanya jika $u\in L^0_J$. Jadi inilah $J^*$
yang kita cari.

\medskip

{\bf (e)} Terapkan kembali (c), kali ini pada $f^{\ssbullet}$.
}%end of proof of 254R

\leader{254S}{Proposisi} Misalkan
$\langle(X_i,\Sigma_i,\mu_i)\rangle_{i\in I}$ suatu keluarga ruang probabilitas,
dengan produk $(X,\Lambda,\lambda)$.

(a) Jika $A\subseteq X$
ditentukan oleh koordinat dalam $I\setminus\{j\}$ untuk setiap $j\in I$,
maka ukuran luarnya $\lambda^*A$ harus bernilai $0$ atau $1$.

(b) Jika $W\in\Lambda$ dan $\lambda W>0$, maka untuk setiap $\epsilon>0$
terdapat $W'\in\Lambda$ dan himpunan berhingga $J\subseteq I$ sedemikian sehingga
$\lambda W'\ge 1-\epsilon$, dan untuk setiap $x\in W'$ terdapat $y\in W$
sedemikian sehingga $x\restr I\setminus J=y\restr I\setminus J$.

\proof{ Untuk $J\subseteq I$, tuliskan $X_J$ untuk $\prod_{i\in J}X_i$ dan
$\lambda_J$ untuk ukuran produk pada $X_J$.

\medskip

{\bf (a)} Misalkan $W$ suatu selubung terukur dari $A$. Menurut 254Rd, terdapat
$J\subseteq I$ terkecil yang untuknya
terdapat $W'\in\Lambda$, yang ditentukan oleh koordinat dalam $J$, dengan
$\lambda(W\symmdiff W')=0$. Sekarang $J=\emptyset$. \Prf\ Ambil sembarang
$j\in I$. Maka $A$ ditentukan oleh koordinat dalam $I\setminus\{j\}$,
yaitu, dapat dipandang sebagai $X_j\times A'$ untuk suatu
$A'\subseteq X_{I\setminus\{j\}}$. Kita juga dapat memandang
$\lambda$ sebagai produk dari $\lambda_{\{j\}}$ dan
$\lambda_{I\setminus\{j\}}$ (254N). Misalkan $\Lambda_{I\setminus\{j\}}$
domain dari $\lambda_{I\setminus\{j\}}$. Menurut 251S,

\Centerline{$\lambda^*A
=\lambda_{\{j\}}^*X_j\cdot\lambda_{I\setminus\{j\}}^*A'
=\lambda_{I\setminus\{j\}}^*A'$.}

\noindent Misalkan $V\in\Lambda_{I\setminus\{j\}}$ suatu selubung terukur dari
$A'$. Maka $W'=X_j\times V$ termasuk dalam
$\Lambda$, memuat $A$, dan berukuran $\lambda^*A$, sehingga
$\lambda(W\cap W')=\lambda W=\lambda W'$ dan $W\symmdiff W'$ bersifat
terabaikan. Pada saat yang sama, $W'$ ditentukan oleh koordinat dalam
$I\setminus\{j\}$. Ini berarti bahwa $J$ harus termuat dalam
$I\setminus\{j\}$. Karena $j$ sembarang, $J=\emptyset$.\ \Qed

Tetapi satu-satunya subhimpunan $X$ yang ditentukan oleh koordinat dalam
$\emptyset$ adalah $X$ dan $\emptyset$. Karena $W$ berbeda dari salah satu
himpunan ini hanya pada suatu himpunan terabaikan, $\lambda^*A=\lambda W\in\{0,1\}$, sebagaimana dinyatakan.

\medskip

{\bf (b)} Tetapkan $\eta=\bover12\min(\epsilon,1)\lambda W$. Menurut 254Fe,
terdapat himpunan terukur $V$, yang ditentukan oleh koordinat dalam suatu subhimpunan
berhingga $J$ dari $I$, sedemikian sehingga $\lambda(W\symmdiff V)\le\eta$. Perhatikan bahwa

\Centerline{$\lambda V\ge\lambda W-\eta\ge\Bover12\lambda W>0$,}

\noindent sehingga

\Centerline{$\lambda(W\symmdiff V)\le\Bover12\epsilon\lambda W
\le\epsilon\lambda V$.}

\noindent Kita dapat mengidentifikasi
$\lambda$ dengan produk c.l.d.\ dari $\lambda_J$ dan
$\lambda_{I\setminus J}$ (254N). Misalkan $\tilde W$,
$\tilde V\subseteq X_J\times X_{I\setminus J}$ himpunan-himpunan yang bersesuaian
dengan $W$, $V\subseteq X$. Maka $\tilde V$ dapat dinyatakan sebagai
$U\times X_{I\setminus J}$ dengan $\lambda_JU=\lambda V>0$. Tetapkan
$U'=\{z:z\in X_{I\setminus J}$, $\lambda_J\tilde W^{-1}[\{z\}]=0\}$.
Maka $U'$ terukur terhadap $\lambda_{I\setminus J}$ (252D(ii) lagi,
karena $\lambda_J$ maupun $\lambda_{I\setminus J}$ lengkap), dan

$$\eqalignno{\lambda_JU\cdot\lambda_{I\setminus J}U'
&\le\int\lambda_J(\tilde W^{-1}[\{z\}]\symmdiff U)
  \lambda_{I\setminus J}(dz)\cr
\displaycause{karena jika $z\in U'$, maka
$\lambda_J(\tilde W^{-1}[\{z\}]\symmdiff U)=\lambda_JU$}
&=\int\lambda_J(\tilde W\symmdiff\tilde V)^{-1}[\{z\}]
  \lambda_{I\setminus J}(dz)\cr
&=(\lambda_J\times\lambda_{I\setminus J})(\tilde W\symmdiff\tilde V)\cr
\displaycause{252D sekali lagi}
&=\lambda(W\symmdiff V)
\le\epsilon\lambda V
=\epsilon\lambda_JU.\cr}$$

\noindent Ini berarti bahwa $\lambda_{I\setminus J}U'\le\epsilon$. Tetapkan
$W'=\{x:x\in X$, $x\restr I\setminus J\notin U'\}$; maka
$\lambda W'\ge 1-\epsilon$. Jika $x\in W'$, maka
$z=x\restr I\setminus J\notin U'$, sehingga $\tilde W^{-1}[\{z\}]$ tidak
kosong, yaitu, terdapat $y\in W$ sedemikian sehingga $y\restr I\setminus J=z$.
Jadi $W'$ ini mempunyai sifat-sifat yang diminta.
}%end of proof of 254S

\leader{254T}{Catatan} \cmmnt{Penting untuk memahami bahwa hasil-hasil
di atas berlaku untuk $L^0$ dan $L^1$ serta
himpunan-himpunan-terukur-hingga-suatu-himpunan-terabaikan, bukan untuk himpunan dan fungsi
itu sendiri. Satu gagasan memang berlaku untuk himpunan dan fungsi, baik
terukur maupun tidak.}

\header{254Ta}{\bf (a)} Misalkan $\langle X_i\rangle_{i\in I}$ suatu keluarga
himpunan dengan produk Kartesius $X$. Untuk setiap $J\subseteq I$, misalkan
$\Cal W_J$
merupakan himpunan semua subhimpunan $X$ yang ditentukan oleh koordinat dalam $J$. Maka
$\Cal W_J\cap\Cal W_K=\Cal W_{J\cap K}$ untuk semua $J$, $K\subseteq I$.
\prooflet{\Prf\ Tentu saja
$\Cal W_J\cap\Cal W_K\supseteq\Cal W_{J\cap K}$, karena
$\Cal W_J\supseteq\Cal W_{J'}$ kapan pun $J'\subseteq J$. Di pihak
lain, andaikan $W\in\Cal W_J\cap\Cal W_K$, $x\in W$, $y\in X$
dan $x\restr J\cap K=y\restr J\cap K$. Tetapkan $z(i)=x(i)$ untuk $i\in J$,
$y(i)$ untuk $i\in I\setminus J$. Maka $z\restr J=x\restr J$, sehingga $z\in
W$. Juga $y\restr K=z\restr K$, sehingga $y\in W$. Karena $x$, $y$
sembarang, $W\in\Cal W_{J\cap K}$; karena $W$ sembarang,
$\Cal W_J\cap\Cal W_K\subseteq\Cal W_{J\cap K}$.\ \QeD}
Dengan demikian, untuk sembarang $W\subseteq X$, $\Cal F=\{J:W\in\Cal W_J\}$ merupakan suatu
filter pada $I$ (kecuali jika $W=X$ atau $W=\emptyset$, dalam hal ini
$\Cal F=\Cal PI$). \cmmnt{Tetapi $\Cal F$ belum tentu mempunyai unsur
terkecil, sebagaimana diperlihatkan oleh contoh berikut.}

\spheader 254Tb Tetapkan $X=\{0,1\}^{\Bbb N}$,

\Centerline{$W=\{x:x\in X$, $\lim_{i\to\infty}x(i)=0\}$.}

\noindent Maka untuk setiap $n\in\Bbb N$, $W$ ditentukan oleh koordinat
dalam
$J_n=\{i:i\ge n\}$. Tetapi $W$ tidak ditentukan oleh koordinat dalam
$\bigcap_{n\in\Bbb N}J_n=\emptyset$. \cmmnt{Perhatikan bahwa

\Centerline{$W=\bigcup_{n\in\Bbb N}\bigcap_{i\ge n}\{x:x(i)=0\}$}

\noindent terukur oleh ukuran biasa pada $X$. Tetapi himpunan ini juga
terabaikan (karena terhitung); dalam 254Rd kita mempunyai $J^*=\emptyset$,
$W'=\emptyset$.}


\leader{*254U}{}\cmmnt{ Sekarang saya dapat menguraikan suatu
contoh tandingan yang menjawab pertanyaan alami yang timbul dari \S251.

\medskip

\noindent}{\bf Contoh} Terdapat ruang ukur yang dapat dilokalkan
$(X,\Sigma,\mu)$ dan suatu
ruang probabilitas $(Y,\Tau,\nu)$ sedemikian sehingga ukuran produk c.l.d.
$\lambda$ pada $X\times Y$ tidak dapat dilokalkan.

\proof{{\bf (a)} Ambil $(X,\Sigma,\mu)$ sebagai ruang dalam 216E, sehingga
$X=\{0,1\}^I$, dengan $I=\Cal PC$ untuk suatu himpunan $C$ yang berkardinalitas lebih besar
daripada $\frak c$. Untuk setiap $\gamma\in C$, tuliskan $E_{\gamma}$ untuk
$\{x:x\in X$, $x(\{\gamma\})=1\}$ (yaitu, $G_{\{\gamma\}}$ dalam
notasi 216Ec); maka $E_{\gamma}\in\Sigma$ dan $\mu E_{\gamma}=1$;
juga setiap himpunan terukur berukuran tak nol beririsan dengan suatu $E_{\gamma}$ dalam
suatu himpunan berukuran tak nol, sedangkan $E_{\gamma}\cap E_{\delta}$
terabaikan untuk semua $\gamma$, $\delta$ yang berbeda (lihat 216Ee).

Misalkan $(Y,\Tau,\nu)$ adalah $\{0,1\}^C$ dengan ukuran biasa (254J). Untuk
$\gamma\in C$, misalkan $F_{\gamma}$ adalah $\{y:y\in Y$, $y(\gamma)=1\}$, sehingga
$\nu F_{\gamma}=\bover12$.
Misalkan $\lambda$ ukuran produk c.l.d.\ pada $X\times Y$, dan
$\Lambda$ domainnya.

\medskip

{\bf (b)} Tinjau keluarga $\Cal W=\{E_{\gamma}\times
F_{\gamma}:\gamma\in C\}\subseteq\Lambda$. \Quer\ Andaikan, jika
mungkin, bahwa $V$ merupakan supremum esensial dari $\Cal W$ dalam $\Lambda$
dalam arti 211G. Untuk $\gamma\in C$, tuliskan
$H_{\gamma}=\{x:V[\{x\}]\symmdiff F_{\gamma}$ terabaikan$\}$.
Karena $F_{\gamma}\symmdiff F_{\delta}$ tidak terabaikan,
$H_{\gamma}\cap H_{\delta}=\emptyset$ untuk semua $\gamma\ne\delta$.

Sekarang $E_{\gamma}\setminus H_{\gamma}$ terabaikan terhadap $\mu$ untuk setiap
$\gamma\in C$. \Prf\
$\lambda((E_{\gamma}\times F_{\gamma})\setminus V)=0$, sehingga
$F_{\gamma}\setminus V[\{x\}]$ terabaikan untuk hampir setiap
$x\in E_{\gamma}$, menurut 252D. Di pihak lain, jika kita menetapkan
$F'_{\gamma}=Y\setminus F_{\gamma}$,
$W_{\gamma}=(X\times Y)\setminus(E_{\gamma}\times F'_{\gamma})$, maka kita
melihat bahwa

\Centerline{$(E_{\gamma}\times F'_{\gamma})\cap(E_{\gamma}\times
F_{\gamma})=\emptyset$,
\quad $E_{\gamma}\times F_{\gamma}\subseteq W_{\gamma}$,}

\Centerline{$\lambda((E_{\delta}\times F_{\delta})\setminus W_{\gamma})
=\lambda((E_{\gamma}\times F'_{\gamma})
  \cap(E_{\delta}\times F_{\delta}))
\le\mu(E_{\gamma}\cap E_{\delta})
=0$}

\noindent untuk setiap $\delta\ne\gamma$, sehingga $W_{\gamma}$ merupakan batas atas
esensial untuk $\Cal W$, dan
$V\cap(E_{\gamma}\times F'_{\gamma})=V\setminus W_{\gamma}$ harus
terabaikan terhadap $\lambda$. Dengan demikian
$V[\{x\}]\setminus F_{\gamma}=V[\{x\}]\cap F'_{\gamma}$
terabaikan terhadap $\nu$ untuk $\mu$-hampir setiap $x\in E_{\gamma}$. Tetapi ini
berarti bahwa $V[\{x\}]\symmdiff F_{\gamma}$ terabaikan terhadap $\nu$ untuk
$\mu$-hampir setiap $x\in E_{\gamma}$, yaitu,
$\mu(E_{\gamma}\setminus H_{\gamma})=0$.\ \Qed

Sekarang tinjau keluarga
$\langle E_{\gamma}\cap H_{\gamma}\rangle_{\gamma\in C}$. Ini merupakan keluarga saling
lepas dari himpunan-himpunan berukuran
hingga dalam $X$. Jika $E\in\Sigma$ berukuran tak nol, terdapat
$\gamma\in C$ sedemikian sehingga
$\mu(E_{\gamma}\cap H_{\gamma}\cap E)=\mu(E_{\gamma}\cap E)>0$. Tetapi
ini berarti bahwa $\Cal E=\{E_{\gamma}\cap H_{\gamma}:\gamma\in C\}$
memenuhi syarat-syarat dalam
213Oa, dan $\mu$ pasti dapat dilokalkan secara ketat; padahal tidak.\ \Bang

\medskip

{\bf (c)} Dengan demikian kita telah menemukan keluarga $\Cal W\subseteq\Lambda$ yang tidak mempunyai
supremum esensial dalam $\Lambda$, dan $\lambda$ tidak dapat dilokalkan.
}%end of proof of 254U

\cmmnt{\medskip

\noindent{\bf Catatan} Jika $(X,\Sigma,\mu)$ dan $(Y,\Tau,\nu)$ adalah sebarang
ruang ukur yang dapat dilokalkan dengan ukuran produk c.l.d.\ yang tidak dapat dilokalkan,
maka versi-versi c.l.d.\ keduanya tetap dapat dilokalkan (213Hb) dan
tetap mempunyai suatu
produk yang tidak dapat dilokalkan (251T), yang tidak mungkin dapat dilokalkan secara ketat;
sehingga setidaknya salah satu faktornya tidak dapat dilokalkan secara ketat (251O).
Jadi setiap
contoh dengan tipe ini harus melibatkan ruang lengkap yang ditentukan secara lokal
dan dapat dilokalkan, tetapi tidak dapat dilokalkan secara ketat, seperti dalam 216E.
}%end of comment

\leader{*254V}{}\cmmnt{ Bersesuaian dengan 251U dan 251Wo, kita mempunyai
hasil berikut tentang pangkat terhitung dari ruang probabilitas tanpa atom.

\medskip

\noindent}{\bf Proposisi} Misalkan $(X,\Sigma,\mu)$ ruang probabilitas tanpa atom
dan $I$ suatu himpunan terhitung. Misalkan $\lambda$ ukuran
probabilitas produk pada $X^I$. Maka
$\{x:x\in X^I$, $x$ injektif$\}$ koterabaikan terhadap $\lambda$.
%used in 495

\proof{ Untuk sembarang pasangan $\{i,j\}$ dari unsur-unsur $I$ yang berbeda, himpunan
$\{z:z\in X^{\{i,j\}}$, $z(i)=z(j)\}$ terabaikan untuk ukuran produk
pada $X^{\{i,j\}}$, menurut 251U. Menurut 254Oa,
$\{x:x\in X^I$, $x(i)=x(j)\}$ terabaikan terhadap $\lambda$. Karena $I$
terhitung, hanya ada terhitung banyaknya pasangan seperti itu $\{i,j\}$, sehingga
$\{x:x\in X^I$, $x(i)=x(j)$ untuk unsur-unsur berbeda $i$, $j\in I\}$
terabaikan, dan komplemennya koterabaikan; tetapi komplemen ini
tepat merupakan himpunan fungsi injektif dari $I$ ke $X$.
}%end of proof of 254V

\exercises{
\leader{254X}{Latihan dasar (a)}
%\spheader 254Xa
Misalkan $\langle(X_i,\Sigma_i,\mu_i)\rangle_{i\in I}$ sebarang keluarga
ruang probabilitas, dengan produk $(X,\Lambda,\lambda)$. Tuliskan $\Cal E$
untuk keluarga subhimpunan $X$ yang dapat dinyatakan sebagai gabungan suatu keluarga berhingga
dari silinder-silinder terukur yang saling lepas. (i) Tunjukkan bahwa jika
$C\subseteq X$
merupakan silinder terukur, maka $X\setminus C\in\Cal E$. (ii) Tunjukkan bahwa
$W\cap V\in\Cal E$ untuk semua $W$, $V\in\Cal E$. (iii) Tunjukkan bahwa
$X\setminus W\in\Cal E$ untuk setiap $W\in\Cal E$. (iv) Tunjukkan bahwa
$\Cal E$ merupakan suatu aljabar subhimpunan $X$. (v) Tunjukkan bahwa untuk sembarang
$W\in\Lambda$, $\epsilon>0$, terdapat $V\in\Cal E$ sedemikian sehingga
$\lambda(W\symmdiff V)\le\epsilon^2$. (vi) Tunjukkan bahwa untuk sembarang
$W\in\Lambda$ dan $\epsilon>0$, terdapat silinder-silinder terukur yang saling lepas
$C_0,\ldots,C_n$ sedemikian sehingga
$\lambda(W\cap C_j)\ge(1-\epsilon)\lambda C_j$ untuk setiap $j$ dan
$\lambda(W\setminus\bigcup_{j\le n}C_j)\le\epsilon$.
\Hint{pilih $C_j$ dari
silinder-silinder terukur yang menyusun suatu himpunan $V$ seperti dalam (v).} (vii) Tunjukkan bahwa
jika $f$, $g$ merupakan fungsi yang dapat diintegralkan terhadap $\lambda$ dan $\int_Cf\le\int_Cg$
untuk setiap silinder terukur $C\subseteq X$, maka $f\leae g$.
\Hint{tunjukkan bahwa $\int_Wf\le\int_Wg$ untuk setiap $W\in\Lambda$.}
%254F

\sqheader 254Xb Misalkan $\langle (X_i,\Sigma_i,\mu_i)$ suatu
keluarga ruang probabilitas, dengan produk $(X,\Lambda,\lambda)$.
Tunjukkan bahwa ukuran luar $\lambda^*$ yang didefinisikan oleh $\lambda$ tepat sama dengan
ukuran luar $\theta$ yang diuraikan dalam 254A, yaitu, bahwa $\theta$ merupakan
ukuran luar reguler.
%254F

\spheader 254Xc Misalkan $\langle (X_i,\Sigma_i,\mu_i)$ suatu
keluarga ruang
probabilitas, dengan produk $(X,\Lambda,\lambda)$. Tuliskan $\lambda_0$ untuk
restriksi $\lambda$ pada $\Tensorhat_{i\in I}\Sigma_i$, dan $\Cal C$
untuk keluarga silinder terukur dalam $X$. Andaikan bahwa
$(Y,\Tau,\nu)$ suatu ruang probabilitas dan $\phi:Y\to X$ suatu fungsi.
(i) Tunjukkan bahwa $\phi$ bersifat \imp\ ketika dipandang sebagai fungsi dari
$(Y,\Tau,\nu)$ ke $(X,\Tensorhat_{i\in I}\Sigma_i,\lambda_0)$ jika dan
hanya jika $\phi^{-1}[C]$ termasuk dalam $\Tau$ dan $\nu\phi^{-1}[C]=\lambda_0C$ untuk
setiap $C\in\Cal C$. (ii) Tunjukkan bahwa $\lambda_0$ adalah satu-satunya ukuran pada
$X$ dengan sifat ini. \Hint{136C.}
%254G

\sqheader 254Xd Misalkan $I$ suatu himpunan dan $(Y,\Tau,\nu)$ ruang
probabilitas lengkap. Tunjukkan bahwa suatu fungsi $\phi:Y\to\{0,1\}^I$ bersifat \imp\
untuk $\nu$ dan ukuran biasa pada $\{0,1\}^I$ jika dan hanya jika
$\nu\{y:\phi(y)(i)=1$ untuk setiap $i\in J\}=2^{-\#(J)}$ untuk setiap $J\subseteq I$
berhingga.
%254J

\sqheader 254Xe Misalkan $I$ sebarang himpunan dan $\lambda$ ukuran biasa pada
$X=\{0,1\}^I$. Definisikan penjumlahan pada $X$ seperti dalam 254Jd. Tunjukkan bahwa
pemetaan $(x,y)\mapsto x+y:X\times X\to X$ bersifat \imp, jika $X\times X$
diberi ukuran produknya.
%254J

\sqheader 254Xf Misalkan $I$ sebarang himpunan dan $\lambda$ ukuran biasa pada
$\Cal PI$. (i) Tunjukkan bahwa pemetaan
$a\mapsto a\symmdiff b:\Cal PI\to\Cal PI$ bersifat \imp\ untuk sembarang
$b\subseteq I$; khususnya, $a\mapsto I\setminus a$ bersifat \imp. (ii)
Tunjukkan bahwa pemetaan
$(a,b)\mapsto a\symmdiff b:\Cal PI\times\Cal PI\to\Cal PI$ bersifat \imp.
%254J

\sqheader 254Xg Tunjukkan bahwa untuk sembarang $q\in[0,1]$ dan sebarang himpunan $I$, terdapat
ukuran $\lambda$ pada $\Cal PI$ sedemikian sehingga
$\lambda\{a:J\subseteq a\}=q^{\#(J)}$ untuk setiap $J\subseteq I$ berhingga.
%254J

\sqheader 254Xh Misalkan $(Y,\Tau,\nu)$ ruang probabilitas
lengkap, dan tuliskan $\mu$ untuk ukuran Lebesgue pada $[0,1]$. Andaikan bahwa
$\phi:Y\to[0,1]$ suatu fungsi sedemikian sehingga $\nu\phi^{-1}[I]$ ada dan
sama dengan $\mu I$ untuk setiap interval $I$ berbentuk
$[2^{-n}k,2^{-n}(k+1)]$, dengan $n\in\Bbb N$ dan $0\le k<2^n$. Tunjukkan
bahwa $\phi$ bersifat \imp\ untuk $\nu$ dan $\mu$.
%254K

\spheader 254Xi Tunjukkan bahwa jika $\tilde\phi:\{0,1\}^{\Bbb N}\to[0,1]$
adalah sebarang bijeksi yang dikonstruksi dengan metode 254K, maka
$\{\tilde\phi^{-1}[E]:E\subseteq[0,1]$ suatu himpunan Borel$\}$ tepat merupakan
aljabar-sigma subhimpunan $\{0,1\}^{\Bbb N}$ yang dibangkitkan oleh himpunan-himpunan
$\{x:x(i)=1\}$ untuk $i\in\Bbb N$.
%254K

\spheader 254Xj Misalkan $\langle X_i\rangle_{i\in I}$ suatu keluarga
himpunan, dan untuk setiap $i\in I$, misalkan $\Sigma_i$ suatu aljabar-sigma
subhimpunan $X_i$. Tunjukkan bahwa untuk setiap
$E\in\Tensorhat_{i\in I}\Sigma_i$, terdapat himpunan terhitung
$J\subseteq I$ sedemikian sehingga $E$
dapat dinyatakan sebagai $\pi_J^{-1}[F]$ untuk suatu $F\in\Tensorhat_{i\in J}\Sigma_i$,
dengan menuliskan $\pi_J(x)=x\restr J\in\prod_{i\in J}X_i$ untuk
$x\in\prod_{i\in I}X_i$.
%254M

\spheader 254Xk(i) Tuliskan $\nu$ untuk ukuran biasa pada
$X=\{0,1\}^{\Bbb N}$. Tunjukkan bahwa untuk sembarang $k\ge 1$, $(X,\nu)$
isomorfik dengan $(X^k,\nu_k)$, dengan $\nu_k$ ukuran pada $X^k$ yang
merupakan ukuran produk yang diperoleh dengan memberi setiap faktor $X$ ukuran
$\nu$. (ii) Dengan menuliskan $\mu_{[0,1]}$ untuk ukuran Lebesgue pada
$[0,1]$, dan seterusnya, tunjukkan bahwa untuk sembarang $k\ge 1$, $([0,1]^k,\mu_{[0,1]^k})$
isomorfik dengan $([0,1],\mu_{[0,1]})$.
%254K, 254N

\spheader 254Xl(i) Dengan menuliskan $\mu_{[0,1]}$ untuk ukuran Lebesgue pada
$[0,1]$, dan seterusnya, tunjukkan bahwa $([0,1],\mu_{[0,1]})$ isomorfik dengan
$(\coint{0,1},\mu_{\coint{0,1}})$. (ii) Tunjukkan bahwa untuk sembarang $k\ge
1$, $(\coint{0,1}^k,\mu_{\coint{0,1}^k})$ isomorfik dengan
$(\coint{0,1},\mu_{\coint{0,1}})$. (iii) Tunjukkan bahwa untuk sembarang $k\ge
1$, $(\Bbb R,\mu_{\Bbb R})$ isomorfik dengan $(\BbbR^k,\mu_{\Bbb
R^k})$.
%254K, 254N

\spheader 254Xm Misalkan $\mu$ ukuran Lebesgue pada $[0,1]$ dan
$\lambda$ ukuran produk pada $[0,1]^{\Bbb N}$. Tunjukkan bahwa
$([0,1],\mu)$ dan $([0,1]^{\Bbb N},\lambda)$ isomorfik.
%254K, 254N

\spheader 254Xn Misalkan $\langle(X_i,\Sigma_i,\mu_i)\rangle_{i\in I}$
suatu keluarga ruang probabilitas lengkap dan $\lambda$ ukuran produk
pada $\prod_{i\in I}X_i$, dengan domain $\Lambda$. Andaikan bahwa
$A_i\subseteq X_i$ untuk
setiap $i\in I$. Tunjukkan bahwa $\prod_{i\in I}A_i\in\Lambda$ jika dan hanya jika (i)
$\prod_{i\in I}\mu_i^*A_i=0$ atau (ii) $A_i\in\Sigma_i$ untuk
setiap $i$ dan $\{i:A_i\ne X_i\}$ terhitung. \Hint{gabungkan
gagasan-gagasan dari 252Xc, 254F, 254L dan 254N.}
%254F, 254L, 254N

\spheader 254Xo Misalkan $\familyiI{(X_i,\Sigma_i,\mu_i)}$ suatu keluarga
ruang probabilitas dengan produk $(X,\Lambda,\lambda)$. (i) Tunjukkan bahwa,
untuk sembarang $A\subseteq X$,

\Centerline{$\lambda^*A=\min\{\lambda_J^*\pi_J[A]:J\subseteq I$
terhitung$\}$,}

\noindent dengan, untuk $J\subseteq I$, saya menuliskan $\lambda_J$ untuk ukuran
probabilitas produk pada $X_J=\prod_{i\in J}X_i$ dan $\pi_J:X\to X_J$ untuk
pemetaan kanonik. (ii) Tunjukkan bahwa jika $J$, $K\subseteq I$ saling lepas
dan $A$, $B\subseteq X$ masing-masing ditentukan oleh koordinat dalam $J$, $K$,
maka $\lambda^*(A\cap B)=\lambda^*A\cdot\lambda^*B$.
%254O

\spheader 254Xp Misalkan $\familyiI{(X_i,\Sigma_i,\mu_i)}$ suatu keluarga
ruang probabilitas dengan produk $(X,\Lambda,\lambda)$. Misalkan $\eusm S$
rentang linear dari himpunan fungsi indikator silinder terukur
dalam $X$, seperti dalam 254Q. Tunjukkan bahwa
$\{f^{\ssbullet}:f\in\eusm S\}$ rapat dalam $L^p(\lambda)$ untuk setiap
$p\in\coint{1,\infty}$.
%254Q

\spheader 254Xq Misalkan $\langle(X_i,\Sigma_i,\mu_i)\rangle_{i\in I}$ suatu
keluarga ruang probabilitas, dan $(X,\Lambda,\lambda)$ produknya;
untuk $J\subseteq I$, misalkan $\Lambda_J$ aljabar-sigma dari anggota-anggota
$\Lambda$ yang ditentukan oleh koordinat dalam $J$, dan
$P_J:L^1=L^1(\lambda)\to L^1_J=L^1(\lambda\restr\Lambda_J)$ ekspektasi
bersyarat yang bersesuaian. (i) Tunjukkan bahwa jika $u\in L^1_J$
dan $v\in L^1_{I\setminus J}$, maka $u\times v\in L^1$ dan
$\int u\times v=\int u\cdot\int v$.
\Hint{253D.} (ii) Tunjukkan bahwa jika $\Cal J\subseteq\Cal PI$ tidak
kosong, dengan $J^*=\bigcap\Cal J$, maka
$L^1_{J^*}=\bigcap_{J\in\Cal J}L^1_J$.
%254R

\spheader 254Xr(i) Misalkan $I$ sebarang himpunan dan $\lambda$ ukuran biasa pada
$\Cal PI$. Misalkan $A\subseteq\Cal PI$ sedemikian sehingga $a\symmdiff b\in A$
kapan pun $a\in A$ dan $b\subseteq I$ berhingga. Tunjukkan bahwa $\lambda^*A$
harus bernilai $0$ atau $1$.
(ii) Misalkan $\lambda$ ukuran biasa pada $\{0,1\}^{\Bbb N}$, dan
$\Lambda$ domainnya. Misalkan $f:\{0,1\}^{\Bbb N}\to\Bbb R$ suatu fungsi
sedemikian sehingga, untuk $x$, $y\in\{0,1\}^{\Bbb N}$,
$f(x)=f(y)\iff\{n:n\in\Bbb N$, $x(n)\ne y(n)\}$ berhingga.
Tunjukkan bahwa $f$ tidak $\Lambda$-terukur.
\Hint{untuk sembarang $q\in\Bbb Q$, $\lambda^*\{x:f(x)\le q\}$ bernilai $0$ atau
$1$.}
%254S

\spheader 254Xs Misalkan $\familyiI{X_i}$ sebarang keluarga himpunan dan
$A\subseteq B\subseteq\prod_{i\in I}X_i$. Andaikan bahwa $A$
ditentukan oleh koordinat dalam $J\subseteq I$ dan bahwa $B$ ditentukan
oleh koordinat dalam $K$. Tunjukkan bahwa terdapat himpunan $C$ sedemikian sehingga
$A\subseteq C\subseteq B$ dan $C$ ditentukan oleh koordinat dalam
$J\cap K$.
%254T

\leader{254Y}{Latihan lanjutan (a)}
%\spheader 254Ya
Misalkan $\familyiI{(X_i,\Sigma_i,\mu_i)}$ suatu keluarga
ruang probabilitas, dan untuk $J\subseteq I$, misalkan $\lambda_J$ ukuran
produk
pada $X_J=\prod_{i\in J}X_i$; tuliskan $X=X_I$, $\lambda=\lambda_I$
dan $\pi_J(x)=x\restr J$ untuk $x\in X$ dan $J\subseteq I$.

\quad (i) Tunjukkan bahwa untuk $K\subseteq J\subseteq I$, kita mempunyai pemetaan alami
yang linear, melestarikan urutan, dan
melestarikan norma, $T_{JK}:L^1(\lambda_K)\to L^1(\lambda_J)$, yang didefinisikan dengan
menuliskan $T_{JK}(f^{\ssbullet})=(f\pi_{KJ})^{\ssbullet}$ untuk setiap fungsi yang
dapat diintegralkan terhadap $\lambda_K$, yaitu $f$, dengan $\pi_{KJ}(y)=y\restr K$ untuk
$y\in X_J$.

\quad (ii) Tuliskan $\Cal K$ untuk himpunan semua subhimpunan berhingga dari $I$. Tunjukkan
bahwa jika $W$ sebarang ruang Banach dan $\langle T_K\rangle_{K\in\Cal K}$
suatu keluarga sedemikian sehingga ($\alpha$) $T_K$ merupakan operator linear terbatas dari
$L^1(\lambda_K)$ ke $W$ untuk setiap $K\in\Cal K$ ($\beta$)
$T_K=T_JT_{JK}$ kapan pun $K\subseteq J\in \Cal K$ ($\gamma$)
$\sup_{K\in\Cal K}\|T_K\|<\infty$, maka terdapat operator linear terbatas yang tunggal
$T:L^1(\lambda)\to W$ sedemikian sehingga $T_K=TT_{IK}$ untuk setiap
$K\in\Cal K$.

\quad (iii) Tuliskan $\Cal J$ untuk himpunan semua subhimpunan terhitung dari $I$.
Tunjukkan bahwa $L^1(\lambda)=\bigcup_{J\in\Cal J}T_{IJ}[L^1(\lambda_J)]$.
%254F

\spheader 254Yb\dvAnew{2010} Misalkan $\familyiI{(X_i,\Sigma_i,\mu_i)}$
sebarang keluarga ruang ukur.
Tetapkan $X=\prod_{i\in I}X_i$ dan misalkan $\Cal F$ suatu filter pada
himpunan $[I]^{<\omega}$ dari subhimpunan berhingga $I$ sedemikian sehingga
$\{J:i\in J\in[I]^{<\omega}\}\in\Cal F$ untuk setiap $i\in I$. Tunjukkan bahwa
terdapat ukuran lengkap yang ditentukan secara lokal
$\lambda$ pada $X$ sedemikian sehingga $\lambda(\prod_{i\in I}E_i)$
terdefinisi dan sama dengan $\lim_{J\to\Cal F}\prod_{i\in J}\mu_iE_i$ kapan pun
$E_i\in\Sigma_i$ untuk setiap $i\in I$ dan
$\lim_{J\to\Cal F}\prod_{i\in J}\mu_iE_i$ terdefinisi dalam
$\coint{0,\infty}$.
\Hint{{\smc Baker 04}.}  %n10313.tex
%254F

\spheader 254Yc
Misalkan $\familyiI{(X_i,\Sigma_i,\mu_i)}$ suatu keluarga ruang probabilitas,
dan $\lambda$ suatu ukuran lengkap pada $X=\prod_{i\in I}X_i$.
Andaikan bahwa untuk setiap ruang probabilitas lengkap $(Y,\Tau,\nu)$ dan
fungsi
$\phi:Y\to X$, $\phi$ bersifat \imp\ untuk $\nu$ dan $\lambda$ jika dan hanya jika
$\nu\phi^{-1}[C]$ terdefinisi dan
sama dengan $\theta_0 C$ untuk setiap silinder terukur $C\subseteq X$,
dengan menuliskan $\theta_0$ untuk fungsional dalam 254A. Tunjukkan bahwa
$\lambda$ merupakan ukuran produk pada $X$.
%254G

\spheader 254Yd Misalkan $I$ suatu himpunan, dan $\lambda$ ukuran biasa
pada $\{0,1\}^I$. Tunjukkan bahwa $L^1(\lambda)$ separabel, dalam
topologi normanya, jika dan hanya jika $I$ terhitung.
%254J

\spheader 254Ye\dvAnew{2013} Misalkan $f:[0,1]\to[0,1]^2$ suatu fungsi
yang bersifat \imp\ untuk ukuran Lebesgue bidang pada $[0,1]^2$ dan
ukuran Lebesgue linear pada $[0,1]$, seperti dalam 134Yl; misalkan $f_1$, $f_2$ adalah
koordinat-koordinat $f$. Definisikan $g:[0,1]\to[0,1]^{\Bbb N}$ dengan menetapkan
$g(t)=\sequencen{f_1f_2^n(t)}$ untuk $0\le t\le 1$. Tunjukkan bahwa $g$ bersifat
\imp. \Hint{tunjukkan bahwa $g_n:[0,1]\to[0,1]^{n+1}$ bersifat \imp\ untuk setiap
$n\ge 1$, dengan
$g_n(t)=(f_1(t),f_1f_2(t),\ldots,\penalty-100f_1f_2^{n-1}(t),f_2^n(t))$
untuk $t\in[0,1]$.}
%254N 254G out of order query

\spheader 254Yf Misalkan $I$ suatu himpunan, dan $\lambda$ ukuran biasa pada
$\Cal PI$. Tunjukkan bahwa jika $\Cal F$ suatu ultrafilter non-prinsipal pada $I$,
maka $\lambda^*\Cal F=1$. \Hint{254Xr, 254Xf.}
%254J, 254S

\spheader 254Yg Misalkan $(X,\Sigma,\mu)$, $(Y,\Tau,\nu)$ dan $\lambda$ seperti
dalam 254U. Tetapkan $A=\{x_{\gamma}:\gamma\in C\}$ sebagaimana didefinisikan dalam 216E. Misalkan
$\mu_A$ ukuran subruang pada $A$, dan $\tilde\lambda$ ukuran produk c.l.d.\
dari $\mu_A$ dan $\nu$ pada $A\times Y$. Tunjukkan bahwa
$\tilde\lambda$ merupakan perluasan sejati dari ukuran subruang
$\lambda_{A\times Y}$. \Hint{tinjau
$\tilde W=\{(x_{\gamma},y):\gamma\in C$, $y\in F_{\gamma}\}$, dalam
notasi 254U.}
%254U

\spheader 254Yh Misalkan $(X,\Sigma,\mu)$ suatu ruang probabilitas tanpa atom,
$I$ suatu himpunan dengan kardinalitas paling besar $\#(X)$, dan $A$ himpunan fungsi
injektif dari $I$ ke $X$. Tunjukkan bahwa $A$ mempunyai ukuran luar penuh untuk
ukuran produk pada $X^I$.
%254V

}%end of exercises

\endnotes{
\Notesheader{254} Walaupun ada banyak alasan untuk mempelajari produk tak hingga
dari ruang probabilitas, dari sudut pandang
teori ukuran abstrak, ada satu alasan yang jauh paling utama: produk-produk itu menyediakan konstruksi
jenis ruang ukur yang pada hakikatnya baru. Saya tidak dapat menguraikan sifat
`kebaruan' ini secara efektif tanpa memasuki wilayah
Jilid 3. Tetapi ruang-ruang fungsi dalam Bab 24 setidaknya menyediakan
suatu ungkapan yang dapat kita gunakan: inilah ruang-ruang {\it probabilitas} pertama
$(X,\Lambda,\lambda)$ yang telah kita lihat yang untuknya $L^1(\lambda)$ belum tentu
separabel dalam topologi normanya (254Yd).

Rumus-rumus dalam 254A, seperti rumus-rumus dalam 251A, mengarah
secara sangat alami pada ukuran; titik ketika
rumus-rumus itu menjadi lebih daripada sekadar keingintahuan adalah ketika kita menemukan bahwa ukuran produk
$\lambda$ merupakan suatu
ukuran probabilitas (254Fa), yang harus dipandang sebagai argumen penting
dalam bagian ini, sebagaimana 251E merupakan dasar utama
\S251. Menurut saya, sungguh luar biasa bahwa tidak ada perbedaan pada
hasil di sini, baik $I$ berhingga, tak hingga terhitung, maupun tak terhitung.
Jika Anda menuliskan bukti untuk kasus $I=\Bbb N$, akan tampak alami
untuk memperluas himpunan-himpunan $J_n$ hingga menjadi segmen awal dari $I$ itu sendiri,
dan dengan demikian menghindari sama sekali himpunan bantu $K$; tetapi ini merupakan
penyederhanaan yang menyesatkan, karena menyembunyikan ciri penting dari
argumen tersebut, yaitu bahwa setiap barisan dalam $\Cal C$ hanya melibatkan terhitung
banyaknya koordinat, sehingga selama kita hanya menangani satu
barisan seperti itu, ketakterhitungan seluruh himpunan $I$ tidak relevan.
Prinsip umum ini secara alami meresapi seluruh bagian; dalam
254O saya telah mencoba menjabarkan bagaimana banyak pertanyaan yang
menarik bagi kita dapat dinyatakan dalam subproduk terhitung dari
ruang-ruang faktor $X_i$. Lihat juga latihan 254Xj, 254Xn, dan
254Ya(iii).

Ada sisi yang agak paradoksal pada prinsip ini: bahkan
subhimpunan-subhimpunan $E_i$ dari $X_i$ yang berperilaku paling baik dapat gagal mempunyai produk terukur
$\prod_{i\in I}E_i$ jika $E_i\ne X_i$ untuk tak terhitung banyaknya $i$. Sebagai
contoh, $\ooint{0,1}^I$ bukan subhimpunan terukur dari $[0,1]^I$ jika $I$
tak terhitung (254Xn). Himpunan ini mempunyai ukuran luar penuh dan
ukuran produknya sendiri tepat merupakan ukuran subruang (254L), tetapi setiap
subhimpunan terukur harus berukuran nol. Intinya adalah bahwa himpunan kosong
merupakan satu-satunya anggota $\Tensorhat_{i\in I}\Sigma_i$, dengan
$\Sigma_i$ aljabar himpunan bagian $[0,1]$ yang terukur Lebesgue untuk
setiap $i$, yang termuat dalam $\ooint{0,1}^I$ (lihat 254Xj).

Seperti dalam \S251, saya menggunakan konstruksi yang secara otomatis menghasilkan
ukuran lengkap pada ruang produk. Saya yakin bahwa inilah pilihan terbaik
untuk `ukuran' produk. Tetapi ada kalanya
restriksinya pada aljabar-sigma yang dibangkitkan oleh silinder-silinder terukur
layak diperhatikan; lihat 254Xc.

Teorema 254G merupakan hasil dengan tipe yang akan lebih sering muncul dalam Jilid 3
daripada dalam jilid ini. Teorema tersebut menguraikan ukuran produk dalam kerangka
bukan mengenai apa ukuran itu {\it adanya}, melainkan apa yang {\it dilakukannya};
secara khusus, dalam kerangka suatu sifat dari keluarga fungsi
\imp\ yang terkait. Oleh
karena itu, teorema ini merupakan suatu `teorema pemetaan universal'. (Bandingkan 253F.)
Karena uraian ini cukup untuk menentukan ukuran produk
sepenuhnya (254Yc), tidaklah mengherankan bahwa saya menggunakannya berulang kali.

`Ukuran biasa' pada $\{0,1\}^I$ (254J) kadang-kadang disebut
`ukuran pelemparan koin' karena dapat digunakan untuk memodelkan konsep
melempar koin berkali-kali secara sembarang dengan indeks himpunan $I$, dengan memandang suatu
$x\in\{0,1\}^I$ sebagai hasil ketika koin menunjukkan `kepala'
tepat untuk semua $i\in I$ yang memenuhi $x(i)=1$. Himpunan-himpunan, atau `kejadian',
dalam kelas $\Cal C$ adalah himpunan yang dapat ditentukan dengan menyatakan
hasil dari
sejumlah berhingga pelemparan, dan probabilitas sebarang barisan tertentu yang terdiri atas
$n$ hasil adalah $1/2^n$, terlepas dari pelemparan mana yang kita perhatikan atau
urutannya. Dalam Bab 27 saya akan kembali menggunakan ukuran produk
untuk merepresentasikan probabilitas yang melibatkan kejadian-kejadian independen.

Dalam 254K saya sampai pada kasus pertama dalam risalah ini mengenai isomorfisme
tak trivial antara dua ruang ukur. Jika Anda dibesarkan dengan
sajian konvensional matematika murni abstrak modern yang berlandaskan aljabar
dan topologi, mungkin Anda sudah memperhatikan ketiadaan
penekanan pada konsep `homomorfisme' atau `isomorfisme'. Di sini
memang saya mulai berbicara tentang `isomorfisme' antara ruang-ruang ukur
bahkan tanpa bersusah payah mendefinisikannya; saya berharap sungguh jelas bahwa
isomorfisme antara ruang ukur $(X,\Sigma,\mu)$ dan
$(Y,\Tau,\nu)$ merupakan bijeksi $\phi:X\to Y$ sedemikian sehingga
$\Tau=\{F:F\subseteq Y$, $\phi^{-1}[F]\in\Sigma\}$ dan
$\nu F=\mu\phi^{-1}[F]$ untuk setiap $F\in\Tau$, sehingga $\Sigma$
tentu sama dengan $\{E:E\subseteq X$, $\phi[E]\in\Tau\}$ dan
$\mu E=\nu\phi[E]$
untuk setiap $E\in\Sigma$. Jika diutarakan seperti ini, Anda mungkin, jika telah mengerjakan
latihan-latihan Jilid 1, teringat pada beberapa konstruksi
aljabar-sigma dalam 111Xc--111Xd dan `ukuran citra' dalam
234C--234D. Hasil dalam 254K (lihat juga 134Yo) secara alami mengarah pada dua
pengertian yang berbeda mengenai
`homomorfisme' antara dua ruang ukur $(X,\Sigma,\mu)$ dan
$(Y,\Tau,\nu)$:

\inset{(i) suatu fungsi $\phi:X\to Y$ sedemikian sehingga $\phi^{-1}[F]\in\Sigma$
dan $\mu\phi^{-1}[F]=\nu F$ untuk setiap $F\in\Tau$,}

\inset{(ii) suatu fungsi $\phi:X\to Y$ sedemikian sehingga $\phi[E]\in \Tau$ dan
$\nu\phi[E]=\mu E$ untuk setiap $E\in\Sigma$.}

\noindent Pada kedua definisi, kita menemukan bahwa bijeksi $\phi:X\to Y$
merupakan isomorfisme jika dan hanya jika $\phi$ dan $\phi^{-1}$ keduanya merupakan homomorfisme.
(Juga, tentu saja, komposisi homomorfisme merupakan suatu
homomorfisme.) Pandangan saya sendiri ialah bahwa (i) lebih penting, dan
dalam risalah ini saya mempelajari fungsi-fungsi tersebut secara panjang lebar, dengan menyebut
fungsi-fungsi itu `\imp'. Tetapi keduanya memiliki kegunaan. Fungsi
$\phi$ dalam 254K bukan hanya memenuhi kedua definisi, melainkan juga
`hampir' merupakan isomorfisme dalam beberapa cara yang berbeda, dan mungkin
yang terpenting ialah bahwa terdapat himpunan-himpunan koterabaikan
$X'\subseteq\{0,1\}^{\Bbb N}$,
$Y'\subseteq [0,1]$ sedemikian sehingga $\phi\restr X'$ merupakan suatu
isomorfisme antara $X'$ dan $Y'$ ketika keduanya
diberi ukuran subruangnya.

Setelah menetapkan isomorfisme antara $[0,1]$ dan
$\{0,1\}^{\Bbb N}$, kita segera memperoleh lebih banyak isomorfisme; lihat
254Xk--254Xm. Bahkan, ukuran Lebesgue pada $[0,1]$ isomorfik dengan
sebagian besar ruang probabilitas yang muncul dalam penerapan. Dalam
Jilid 3 dan 4 saya akan membahas isomorfisme-isomorfisme ini secara panjang lebar.

Gagasan umum `subproduk' berhubungan dengan beberapa hasil
yang paling mendalam dan paling khas dalam teori ukuran produk.
Karena kita memperhatikan produk dari sebarang keluarga
ruang probabilitas, definisinya harus mengabaikan setiap struktur yang mungkin ada dalam
himpunan indeks $I$ dari 254A--254C. Tetapi banyak penerapan, secara cukup
alami, menangani himpunan indeks yang mempunyai subhimpunan atau partisi yang diistimewakan, dan
langkah penting pertama adalah `hukum asosiatif' (254N; bandingkan
251Xe--251Xf dan 251Wh). Misalnya, inilah alat yang memungkinkan
kita menerapkan teorema Fubini di dalam produk tak hingga. Pemetaan-pemetaan
proyeksi alami dari $\prod_{i\in I}X_i$ ke $\prod_{i\in J}X_i$, dengan
$J\subseteq I$, saling berhubungan dengan cara yang telah digunakan sebagai
dasar teorema-teorema dalam \S235; ukuran produk pada $\prod_{i\in J}X_i$
tepat merupakan bayangan dari ukuran produk pada $\prod_{i\in I}X_i$
(254Oa). Dalam 254O--254Q saya menyelidiki konsekuensi dari fakta ini dan
fakta yang telah dicatat bahwa semua himpunan terukur dalam produk
`pada hakikatnya' ditentukan oleh koordinat dalam suatu himpunan terhitung.

Dalam 254R saya menyelidiki lebih jauh gagasan tentang suatu himpunan
$W\subseteq\prod_{i\in I}X_i$ yang `ditentukan oleh koordinat dalam' suatu himpunan
$J\subseteq I$. Dalam bentuk primitifnya, ini merupakan gagasan teori-himpunan murni
(254M, 254Ta). Saya rasa bahkan himpunan $I$ yang berunsur tiga dapat
memberi kita kejutan; saya mengajak Anda mencoba memvisualisasikan subhimpunan
$[0,1]^3$ yang ditentukan oleh pasangan-pasangan koordinat. Tetapi
interaksi gagasan ini dengan gagasan teori-ukur, dan khususnya
dengan kesediaan menambah atau membuang himpunan terabaikan, menghasilkan jauh
lebih banyak, dan khususnya himpunan koordinat minimal yang tunggal
yang terkait dengan himpunan dan fungsi terukur (254R). Tentu saja
hasil-hasil ini dapat diuraikan secara anggun dan efektif dalam kerangka
ruang $L^1$ dan $L^0$, tempat himpunan-himpunan terabaikan disingkirkan dari pandangan
ketika ruang-ruang itu dikonstruksi. Dasar semua ini ialah fakta bahwa
operator-operator ekspektasi bersyarat yang terkait dengan subproduk
saling dikalikan dengan cara yang sesederhana mungkin (254Ra); tetapi diperlukan
gagasan lebih lanjut untuk menunjukkan bahwa jika $\Cal J$ suatu keluarga tak kosong dari
subhimpunan $I$, maka $L^0_{\bigcap\Cal J}=\bigcap_{J\in\Cal J}L^0_J$
(lihat bagian (b) dari bukti 254R; sumber juga mengacu ke 254Xq(iii),
tetapi bagian tersebut tidak ada dalam seksi yang dirujuk).

254Sa merupakan suatu versi dari `hukum nol-satu' (272O di bawah). 254Sb merupakan
versi kuat dari prinsip bahwa himpunan terukur dalam suatu
produk harus dapat dihampiri oleh himpunan yang ditentukan oleh suatu himpunan koordinat
{\it berhingga} (254Fe, 254Qa, 254Xa). Jelas bukan suatu kebetulan
bahwa himpunan $W$ dalam 254Tb terabaikan. Dalam \S272 saya akan meninjau kembali banyak
gagasan dalam 254R--254S dan 254Xq, khususnya, dalam konteks yang lebih
umum berupa `aljabar-sigma independen'.

Terakhir, 254U dan 254Yg hampir sama sekali tidak termasuk dalam bagian ini; keduanya merupakan
persoalan yang belum selesai dari \S251. Keduanya ditempatkan di sini karena konstruksi
dalam 254A--254C merupakan cara paling sederhana untuk menghasilkan ruang
probabilitas $(Y,\Tau,\nu)$ yang cukup kompleks.
}%end of notes

\discrpage
